\chapter{Thiết kế MAC-F}
\label{ch:design}

\section{Đường dữ liệu và giao thức}
Hình~\ref{fig:mac-f-datapath} tách rạch ròi đường dữ liệu, metadata và điều
khiển, sử dụng hộp viền đôi để phân biệt các khối con đường ống nhiều tầng
(\texttt{Parser P1/P2}, \texttt{FULL\_TERM T1--T3}, \texttt{posit\_add S1--S5},
\texttt{posit\_pack P1/P2}) với mạch tổ hợp thuần (viền đơn) và thanh ghi D-FF
(ô xanh đậm), kèm thanh ranh giới chu kỳ $14\text{--}25$ và chú giải ký hiệu
chung ở cuối hình. Các khối, độ rộng và vị trí thanh ghi dưới đây được trích
trực tiếp từ RTL của đề tài; phương pháp pipeline bộ giải mã--chuẩn hóa Posit
trên FPGA kế thừa cấu trúc đã được khảo sát trong các công trình của
Jaiswal--So, Podobas--Matsuoka và FloPoCo \cite{jaiswal2018posit,podobas2018posit,murillo2020flopoco}.
\rtlfig{Datapath kiến trúc chính MAC-F. Khối viền đôi là mô-đun con đường ống nhiều tầng nội bộ (gồm \sig{Parser} P1/P2, \sig{FULL\_TERM} T1--T3, \sig{posit\_add} S1--S5 và \sig{posit\_pack} P1/P2); ô xanh đậm là thanh ghi; nét liền mang dữ liệu, nét chấm mang metadata và nét đứt mang điều khiển. Hình~\ref{fig:nf-f-split} phóng to chỗ khác nhau giữa NF và F.}{fig:mac-f-datapath}{
\node[pin] (a) at (-4.7,0) {\sig{a[31:0]}};
\node[pin] (b) at (-.3,0) {\sig{b[31:0]}};
\node[pin] (c) at (4.7,0) {\sig{c[31:0]}};
\node[pipe,text width=32mm] (pa) at (-4.7,-1.05) {\hw{Parser A}\\P1/P2 (2 tầng)};
\node[pipe,text width=32mm] (pb) at (-.3,-1.05) {\hw{Parser B}\\P1/P2 (2 tầng)};
\node[pipe,text width=32mm] (pc) at (4.7,-1.05) {\hw{Parser C}\\P1/P2 (2 tầng)};
\foreach \n in {a,b,c}{\draw[wire] (\n)--(p\n);}
\node[unit,text width=54mm] (ops) at (-2.5,-3) {\hw{OPS: chọn vai X/Y}\\$s_A\oplus s_B,\quad sf_A+sf_B$};
\node[ff,text width=35mm] (ctx) at (4.7,-2.85) {\hw{Context C}\\frac 26, $sf$:11\\$s,Z,N$};
\draw[wire] (pa.south) -- ++(0,-.53) -| (ops.150) node[pos=.22,tag,above]{12, $s,sf$};
\draw[wire] (pb.south) -- ++(0,-.53) -| (ops.30) node[pos=.22,tag,above]{12, $s,sf$};
\draw[wire] (pc)--node[tag,right]{\sig{\{frac,frac\_lo\}}}(ctx);
\node[unit,text width=58mm] (mx) at (-2.5,-4.7) {\hw{MUX Zero/NaR: chọn X,Y hoặc 0}};
\draw[wire] (ops)--node[tag,right]{X:12, Y:13}(mx);
\node[unit,text width=32mm] (exc) at (4.7,-4.7) {\hw{Ngoại lệ từ Parser}\\$Z=Z_A\lor Z_B$\\$N=N_A\lor N_B\lor N_C$};
\draw[enable] (exc.west)--(mx.east);
\node[pipe,text width=58mm] (mul) at (-2.5,-6.35) {\hw{Nhân lặp FULL\_TERM} (T1--T3)\\SAC $\to$ SBM $\to$ cộng Q24};
\draw[wire] (mx)--(mul);
\node[unit,text width=58mm] (norm) at (-2.5,-8.05) {\hw{Adapter tích}\\chọn scale và fraction theo \sig{acc[25]}};
\draw[wire] (mul)--node[tag,right]{\sig{acc[25:0]}, $sf_{\mathrm{core}}$}(norm);
\node[ff,text width=34mm] (pm) at (4.7,-8.05) {\hw{Metadata tích}\\$Z,N$; $sf_{\mathrm{core}}+1$};
\draw[meta] (exc.south)--(pm.north);
\draw[meta] (pm.west)--node[tag,above]{$sf_{\mathrm{core}}+1$}(norm.east);
\node[pipe,text width=76mm] (add) at (0,-10.06) {\hw{Bộ cộng}~\sig{posit\_add}~\hw{(5 tầng: S1--S5)}\\tích + C: căn thang, cộng/trừ, chuẩn hóa};
\draw[wire] (norm.south)--node[tag,right]{frac 26, $sf$:11}(add.north -| norm);
\coordinate (cport) at (2.7,0);
\coordinate (zport) at (.62,0);
\draw[wire] (ctx.east) -- (7.22,-2.85) -- (7.22,-9.17) -- node[tag,below]{C giữ đến S1} (2.7,-9.17) -- (add.north -| cport);
\draw[meta] (pm.south) -- (4.7,-8.83) -- node[tag,above]{$Z,N$} (.62,-8.83) -- (add.north -| zport);
\node[pipe,text width=76mm] (pack) at (0,-11.94) {\hw{Packer}~\sig{posit\_pack}~\hw{(2 tầng: P1/P2)}\\RNE và chọn mã ngoại lệ};
\draw[wire] (add)--node[tag,right]{frac 53, $sf$:11, sticky, cờ}(pack);
\node[pin] (out) at (0,-13.06) {\sig{d[31:0], flags[4:0]}};
\draw[wire] (pack)--(out);
\node[explain,text width=29mm] (ctl) at (-5.87,-10.4) {\hw{Wrapper}\\\sig{pair\_accept}: chốt cfg\\\sig{pair\_retire}: chốt C, launch\\\sig{retire}: nhả core};
\draw[enable] (ctl.north)|-(mul.west);
\node[explain,text width=29mm] at (5.89,-11.94) {\hw{Biên đầu ra}\\\sig{out\_valid} giữ cùng D/cờ khi \sig{out\_ready=0}};
\draw[draw=dpline,fill=dpcomb!50,rounded corners=1.5pt] (-7.4,-13.48) rectangle (7.4,-14.08);
\node[anchor=west,font=\fontsize{7.8}{9}\selectfont\sffamily\bfseries,text=dpink] at (-7.25,-13.78) {Chu kỳ (14--25):};
\draw[dpline] (-3.85,-13.52)--(-3.85,-14.04);
\node[anchor=center,font=\fontsize{7.8}{9}\selectfont\sffamily,text=dpink] at (-2.78,-13.78) {\sig{P1|P2} (2)};
\draw[dpline] (-1.70,-13.52)--(-1.70,-14.04);
\node[anchor=center,font=\fontsize{7.8}{9}\selectfont\sffamily,text=dpink] at (-0.05,-13.78) {\sig{T1..Iter(t)..T3} (3+max(1,t))};
\draw[dpline] (1.60,-13.52)--(1.60,-14.04);
\node[anchor=center,font=\fontsize{7.8}{9}\selectfont\sffamily,text=dpink] at (2.68,-13.78) {\sig{S1..S5} (5)};
\draw[dpline] (3.75,-13.52)--(3.75,-14.04);
\node[anchor=center,font=\fontsize{7.8}{9}\selectfont\sffamily,text=dpink] at (4.78,-13.78) {\sig{P1|P2} (2)};
\draw[dpline] (5.80,-13.52)--(5.80,-14.04);
\node[anchor=center,font=\fontsize{7.8}{9}\selectfont\sffamily,text=dpink] at (6.60,-13.78) {Retire (1)};
\draw[draw=dpline,rounded corners=2pt,fill=white] (-7.4,-14.22) rectangle (7.4,-15.72);
\node[anchor=west,font=\fontsize{8.2}{9.6}\selectfont\sffamily\bfseries,text=dpink] at (-7.2,-14.50) {Chú giải ký hiệu (chung cho các sơ đồ datapath):};
\node[unit,minimum height=4.0mm,minimum width=6.5mm,inner sep=1pt] at (-6.7,-14.92) {};
\node[anchor=west,font=\fontsize{8.0}{9.4}\selectfont\sffamily] at (-6.22,-14.92) {Tổ hợp};
\node[pipe,minimum height=4.0mm,minimum width=6.5mm,inner sep=1pt] at (-4.55,-14.92) {};
\node[anchor=west,font=\fontsize{8.0}{9.4}\selectfont\sffamily] at (-4.07,-14.92) {Đường ống đa tầng};
\node[ff,minimum height=4.0mm,minimum width=6.5mm,inner sep=1pt] at (-1.05,-14.92) {};
\node[anchor=west,font=\fontsize{8.0}{9.4}\selectfont\sffamily] at (-0.57,-14.92) {Thanh ghi D-FF};
\node[pick,minimum height=4.0mm,minimum width=6.5mm,inner sep=1pt] at (2.05,-14.92) {};
\node[anchor=west,font=\fontsize{8.0}{9.4}\selectfont\sffamily] at (2.53,-14.92) {MUX};
\node[explain,minimum height=4.0mm,minimum width=6.5mm,inner sep=1pt] at (3.95,-14.92) {};
\node[anchor=west,font=\fontsize{8.0}{9.4}\selectfont\sffamily] at (4.43,-14.92) {Ghi chú / giao diện};
\draw[wire] (-7.05,-15.40) -- (-6.35,-15.40);
\node[anchor=west,font=\fontsize{8.0}{9.4}\selectfont\sffamily] at (-6.22,-15.40) {Luồng dữ liệu chính};
\draw[meta] (-2.90,-15.40) -- (-2.20,-15.40);
\node[anchor=west,font=\fontsize{8.0}{9.4}\selectfont\sffamily] at (-2.07,-15.40) {Metadata / cờ ($Z,N$)};
\draw[enable] (1.35,-15.40) -- (2.05,-15.40);
\node[anchor=west,font=\fontsize{8.0}{9.4}\selectfont\sffamily] at (2.18,-15.40) {Tín hiệu \sig{enable} / điều khiển};
}


Wrapper nhận một giao dịch khi cả ba Parser và đường xử lý đã sẵn sàng.
Các Parser nhận A/B/C cùng lúc. Wrapper chốt \texttt{cfg\_n} và \texttt{cfg\_ops}
ở cạnh handshake. Khi cả ba Parser hợp lệ và bộ nhân nhận được dữ liệu,
wrapper chuyển A/B sang bộ nhân và giữ dấu, số mũ, fraction và cờ của C
để dùng khi tích sẵn sàng. NaR của bất kỳ toán hạng nào và Zero ở A/B
ép X/Y bằng 0 ở ngõ vào core; core vẫn phát token init-only.
Wrapper chuyển cờ ngoại lệ cùng kết quả vào bộ cộng.

Tại \texttt{core\_done}, tích nội bộ được chuẩn hóa và đưa cùng C vào bộ cộng.
Packer nhận kết quả S5 và giữ mã kết quả, cờ và \texttt{out\_valid} khi
\texttt{out\_ready=0}. Wrapper giải phóng giao dịch khi response retire.
Reset xóa các valid và context đang chờ, sau đó các tầng mở lại qua
các bộ đồng bộ reset và tín hiệu sẵn sàng. Cấu hình OPS ngoài 0/1 hoặc
n ngoài 0--12 bị chặn ở đầu vào.

\subsection{Các ranh giới giữ context và lựa chọn kiến trúc single-outstanding}
\label{subsec:single-outstanding}

Các sơ đồ trong chương này mô tả trực tiếp mã nguồn RTL hoàn chỉnh của lõi
\texttt{MAC-F} cùng hai kiến trúc đối chứng. Cần phân biệt rõ giữa khả năng co
giãn ở cấp mô-đun con và lịch điều phối ở cấp wrapper IP: ba khối chức năng
dùng chung gồm bộ giải mã \texttt{posit\_parser\_mul} ($2$ tầng), bộ cộng
\texttt{posit\_add} ($5$ tầng) và bộ đóng gói \texttt{posit\_pack} ($2$ tầng)
được thiết kế chuẩn hóa theo giao thức bắt tay \texttt{valid/ready} co giãn đầy
đủ (\textit{elastic}) nhằm hỗ trợ ghép nối linh hoạt theo nhịp $II=1$ ở từng
tầng đơn lẻ. Tuy nhiên, ở cấp tích hợp IP \texttt{MAC-NF/F/Quire}, toàn lõi sử
dụng nhân lặp dịch--cộng tuần tự \texttt{mul\_iter\_core} ($1\text{--}12$ chu kỳ
phi đường ống) để giữ \texttt{0 DSP}, đồng thời các bài toán tích lũy chuỗi
($S_i = S_{i-1} + A_i B_i$) có phụ thuộc dữ liệu mang vòng lặp nên mỗi giao dịch
kế tiếp phụ thuộc trực tiếp kết quả $D_{i-1}$. Vì hai lý do này, wrapper chủ động
áp đặt cơ chế \textit{single-outstanding} (chỉ phát một giao dịch đang xử lý tại
một thời điểm) để tối ưu tiết kiệm logic đệm và bảo vệ ngữ cảnh $C$, chấp nhận giữ
các nhánh \textit{backpressure} liên tầng của các khối con ở trạng thái chờ tĩnh.
Độ sâu pipeline của các khối cộng và đóng gói vì vậy quyết định trực tiếp độ
trễ đơn giao dịch, không suy ra $II=1$ cho toàn IP.

\rtlfig{Ba mốc điều khiển trong \sig{mac\_f\_wrapper}. \sig{pair\_retire} chuyển payload từ Parser vào core và context C; \sig{retire} giải phóng giao dịch sau Packer.}{fig:context-control}{
\node[unit,text width=38mm] (accept) at (-5,0) {\sig{pair\_accept}\\\sig{in\_valid} $\land$ \sig{in\_ready}};
\node[unit,text width=51mm] (launch) at (0,0) {\sig{pair\_retire}\\\sig{valid\_a} $\land$ \sig{valid\_b} $\land$ \sig{valid\_c}\\$\land$ \sig{core\_ready}};
\node[unit,text width=38mm] (ret) at (5,0) {\sig{retire}\\\sig{out\_valid} $\land$ \sig{out\_ready}};
\node[ff,text width=38mm] (cfg) at (-5,-1.58) {\hw{Context cấu hình}\\\sig{n\_q[3:0]}, \sig{ops\_q}};
\node[ff,text width=51mm] (ctx) at (0,-1.58) {\hw{Context C và tích}\\frac C: 26; $sf_C$: 11; $s,Z,N$\\cờ tích, $sf_{\mathrm{core}}+1$};
\node[unit,text width=38mm] (core) at (5,-1.58) {\sig{mul\_iter\_core}\\launch, release\_result};
\draw[enable] (accept)--(cfg);
\draw[enable] (launch)--(ctx);
\draw[enable] (ret)--(core);
\draw[enable] (launch.east)--(2.83,0)|-(core.west);
\node[ff,text width=126mm] (busy) at (0,-3.15) {\sig{busy\_q} \quad
 \sig{pair\_accept}: đặt 1 \qquad|\qquad \sig{retire} hoặc reset: xóa};
\draw[enable] (accept.west)--++(-.29,0)|-(busy.west);
\draw[enable] (ret.east)--++(.29,0)|-(busy.east);
\node[tag] (gate) at (0,-4.25) {\sig{in\_ready} yêu cầu \sig{!busy\_q}, cfg hợp lệ, reset đã nhả và các khối sẵn sàng};
\draw[enable] (busy)--(gate);
}



Đầu vào cấu hình được chốt ở cạnh \texttt{pair\_accept}; các đầu vào A/B/C
vào ba Parser cùng cạnh này. \texttt{pair\_retire} là handshake của cụm
Parser với core, không phải retire ở giao diện ngoài. Ở cạnh đó, wrapper
giữ $sf_C[10:0]$, fraction $[25:0]$, dấu và metadata của C. Core nhận
X/Y đã qua OPS hoặc token bypass. Sau khi core phát \texttt{core\_done},
adapter đưa tích và context C vào S1. Wrapper chỉ nhả slot bận khi
\texttt{out\_valid} và \texttt{out\_ready} cùng bằng 1.

\begin{engineeringtable}{Chuyển quyền sở hữu dữ liệu ở các ranh giới của MAC-F.}{tab:mac-f-ownership}
\begin{tabularx}{\linewidth}{@{}>{\raggedright\arraybackslash}p{34mm}Y Y@{}}
\toprule
Sự kiện & Payload được chốt & Điều kiện giữ và giải phóng \\
\midrule
\texttt{pair\_accept} & Mã A/B/C trong Parser, n 4 bit và OPS 1 bit hợp lệ trong wrapper. & Slot bận tới retire hoặc reset. \\
\texttt{pair\_retire} & A/B đã giảm fraction, C đủ 26 bit, tổng scale và metadata đặc biệt. & C ổn định tới lúc S1 nhận tích. \\
\texttt{core\_done} & Accumulator Q24, số mũ và cờ được chuyển cùng C vào adder. & Assertion bảo vệ điều kiện nhận của adder. \\
\texttt{out\_valid} & Mã $D$ và flags 5 bit. & Giữ ổn định khi \texttt{out\_ready=0}, retire đúng một lần khi ready. \\
\bottomrule
\end{tabularx}
\end{engineeringtable}

\subsection{Backpressure, hủy giao dịch và trở lại idle}

Một giao dịch hợp lệ bắt đầu tại cạnh nhận. Nếu sink kéo ready thấp
sau đó, driver không tái cấp đầu vào đang xử lý. Trong trường hợp kết
quả đã đến Packer P2, P2 không ghi đè slot đang giữ. \texttt{out\_valid}
và payload được kiểm cùng nhau, vì giữ đúng mã nhưng hạ valid trước
handshake vẫn làm mất giao dịch. Nếu reset được assert, giao dịch
chưa retire bị hủy. Sau khi mở reset, các đồng bộ nội bộ và điều kiện
sẵn sàng chặn nhận cho đến khi pipeline có thể tiếp tục.

Trong kịch bản kiểm chứng giao thức trên bộ dữ liệu mẫu \texttt{FT-05}, phía nhận giữ \texttt{out\_ready=0} khi
kết quả xuất hiện. Ở các cạnh 14, 15 và 16 tính từ accept, \texttt{d} giữ nguyên
\texttt{1ae34800} với flags \texttt{00}. Cạnh 17 retire đúng kết quả
đó. Ca reset riêng hạ reset sau launch khi token chưa ra core,
sau đó kiểm 30 cạnh không có response cũ. Lượt sau reset dùng lại
cặp FT-05 và trả đúng mã. Trace này minh họa trực tiếp một kịch bản bắt tay
tiêu biểu ở phần thiết kế, còn thống kê diện rộng được trình bày ở Chương~\ref{ch:evaluation}.


\section{Giải mã Posit}
\begin{reportpseudocode}{alg:parser}{Giải mã Posit(32,3) hai tầng (Two-Stage Posit(32,3) Decoding)}
  {Từ mã Posit $p[31:0]$.}
  {Dấu $s$, hệ số tỉ lệ $sf$, hai lát phần phân số $f_H[11:0]$, $f_L[13:0]$ và các cờ $z$, $\eta$, $\mathrm{cut}$.}
  \begin{tabularx}{\linewidth}{@{}r >{\raggedright\arraybackslash}X >{\raggedright\arraybackslash}p{.33\linewidth}@{}}
    \algline{1}{\textbf{P1} \quad $s\gets p[31]$, $u\gets (p[30:0]=0)$}{Tách bit dấu và kiểm tra tải dữ liệu.}
    \algline{2}{$z\gets \neg s\land u$, $\eta\gets s\land u$}{Phát hiện các trường hợp Zero và NaR.}
    \algline{3}{$m\gets p$ \textbf{if} $s=0$ \textbf{else} $(-p\bmod 2^{32})$}{Lấy trị tuyệt đối không dấu 32 bit.}
    \algline{4}{$r\gets m[30]$, $v\gets m[28:0]$}{Đọc bit regime và hạt dịch trái 29 bit.}
    \algline{5}{$d\gets (m[29:0]\oplus m[30:1])\mathbin{\|}1$}{Đánh dấu biên chuyển đổi regime và thêm bit chốt.}
    \algline{6}{$c\gets \mathrm{LZC}(d)$}{Đếm số bit 0 dẫn đầu; chốt tầng P1.}
    \algline{7}{\textbf{P2} \quad $w[28:0]\gets v\ll c$, $e\gets w[28:26]$}{Dịch chuẩn hóa và trích 3 bit mũ.}
    \algline{8}{$k[5:0]\gets \mathop{\mathrm{NOT}}\bigl(\{0,c\}\oplus\{6\{r\}\}\bigr)$}{Tính $k$ không rẽ nhánh ($c$ nếu $r=1$, $-c-1$ nếu $r=0$).}
    \algline{9}{\textbf{if} $z\lor\eta$ \textbf{then} $sf\gets 0$ \textbf{else} $sf\gets (k\ll 3)+e$}{Ghép hệ số tỉ lệ có dấu khi toán hạng hữu hạn.}
    \algline{10}{$f_H\gets w[25:14]$, $f_L\gets w[13:0]$}{Chia phần phân số thành lát cao và lát thấp.}
    \algline{11}{$\mathrm{cut}\gets(f_L\ne 0)$}{Bật cờ cắt đuôi nếu lát thấp khác 0.}
    \algline{12}{\textbf{return} $s,sf,f_H,f_L,z,\eta,\mathrm{cut}$}{Chốt ngõ ra tại tầng đường ống P2.}
  \end{tabularx}
\end{reportpseudocode}
Đường A/B dùng $f_H$ và \texttt{tail\_cut}. Đường C ghép $f_H$ với $f_L$.
Các cờ Zero/NaR điều khiển cách wrapper sử dụng payload.
P1 và P2 giữ metadata đồng bộ với dữ liệu khi tầng sau chưa sẵn sàng.

\subsection{P1 lấy độ lớn và tìm chuyển tiếp}

Kiến trúc Parser dùng chung cấu trúc \texttt{posit\_parser\_mul} cho cả ba toán hạng A, B và C.
Ở P1, dấu và hai cờ đặc biệt đến từ mã gốc. Độ lớn được tạo bằng
XOR với mặt nạ dấu và cộng bù, dùng đường carry của Artix-7.
Chuỗi XOR giữa hai bit kề nhau đánh dấu vị trí thay đổi từ run
đầu sang terminator. Vì vậy cùng một detector tìm chuyển tiếp
cho cả run 0 và run 1. Thêm sentinel 1 ở cuối giúp định nghĩa
được kết quả khi mã độ lớn không có terminator trong 31 bit.

\rtlfig{Hai tầng của \sig{posit\_parser\_mul}. Vector chuyển tiếp chỉ cấp LOD; seed để dịch giữ nguyên từ magnitude. Quire dùng \sig{posit\_parser} với cùng phép dò regime và trả fraction 26 bit.}{fig:parser-stages}{
\node[pin] (p) at (0,.42) {\sig{p[31:0]}};
\node[unit,text width=62mm] (mag) at (-3.05,-1.1) {\hw{Lấy độ lớn}\\$m=(p\oplus\{32\{s\}\})+s$};
\node[unit,text width=52mm] (exc) at (4.2,-1.1) {\hw{Dấu và đặc biệt}\\$s=p[31]$, $u=\lnot (|p[30:0])$\\$Z=\lnot s\land u$, $N=s\land u$};
\draw[wire] (p.south)--++(0,-.2) -| (mag.north);
\draw[wire] (p.south)--++(0,-.2) -| (exc.north);
\node[unit,text width=62mm] (xor) at (-3.05,-2.95) {\hw{Đánh dấu chuyển tiếp}\\$d=\{m[29:0]\oplus m[30:1],1'b1\}$};
\draw[wire] (mag)--node[tag,right]{$m[31:0]$}(xor);
\node[unit,text width=62mm] (lod) at (-3.05,-4.47) {\hw{LOD regime 31 bit}\\\sig{posit\_parser\_f\_regime\_lod}};
\draw[wire] (xor)--node[tag,right]{$d[30:0]$}(lod);
\node[ff,text width=125mm] (p1) at (0,-6.18) {\hw{P1} \quad \sig{cnt[4:0]} \enspace|\enspace
 $r=m[30]$ \enspace|\enspace \sig{seed[28:0]} \enspace|\enspace $s,Z,N$};
\draw[wire] (lod)--node[tag,right]{\sig{cnt[4:0]}}(p1.north -| lod);
\coordinate (seedport) at (.85,0);
\draw[wire] (mag.east) -- (.85,-1.1) -- (p1.north -| seedport)
    node[pos=.55,tag,right]{$m[30], m[28:0]$};
\draw[meta] (exc.south)--(p1.north -| exc) node[pos=.5,tag,right]{$s,Z,N$};
\node[unit,text width=50mm] (sh) at (-3.66,-7.84) {\hw{Dịch trái 29 bit}\\$w=\sig{seed}\ll\sig{cnt}$};
\node[unit,text width=52mm] (reg) at (3.66,-7.84) {\hw{Tái tạo regime signed}\\$r=1$: $k=\sig{cnt}$\\$r=0$: $k=-\sig{cnt}-1$};
\draw[wire] (p1.south -| sh)--(sh);
\draw[wire] (p1.south -| reg)--(reg);
\node[unit,text width=53mm] (slice) at (-3.66,-9.73) {\hw{Tách fraction}\\$f_H=w[25:14]$, $f_L=w[13:0]$\\\sig{tail\_cut} $=|f_L$};
\node[unit,text width=52mm] (scale) at (3.66,-9.73) {\hw{Ghép scale}\\$sf=8k+e$, $e=w[28:26]$\\chọn $k=0$ nếu $Z\lor N$};
\draw[wire] (sh)--(slice);
\draw[wire] (reg)--(scale);
\draw[wire] (sh.east)--(0,-7.84) |- (scale.west) node[pos=.49,tag,right]{$e$:3};
\node[ff,text width=125mm] (p2) at (0,-11.49) {\hw{P2} \quad \sig{sf[10:0]} \enspace|\enspace
 \sig{frac[11:0]}, \sig{frac\_lo[13:0]} \enspace|\enspace \sig{tail\_cut}, $s,Z,N$};
\draw[wire] (slice)--(p2.north -| slice);
\draw[wire] (scale)--(p2.north -| scale);
\draw[meta] (p1.east)--(6.91,-6.18)--node[tag,left]{$s,Z,N$}(6.91,-11.49)--(p2.east);
\node[pin] (ab) at (-3.66,-12.81) {A/B: $f_H$, \sig{tail\_cut}};
\node[pin] (c) at (3.66,-12.81) {C: \sig{\{frac,frac\_lo\}}};
\draw[wire] (p2.south -| ab)--(ab);
\draw[wire] (p2.south -| c)--(c);
}



P1 chốt chỉ số \texttt{cnt}, bit regime đầu, seed \texttt{mag[28:0]}
và dấu/cờ. Seed lấy từ độ lớn gốc. Vector XOR dùng riêng để dò
vị trí và không đi vào đường fraction. Đây là điểm quan trọng nếu
đối chiếu hình với RTL: dữ liệu dịch ở P2 là seed của magnitude,
không phải kết quả XOR. Với mẫu gần cực trị, sentinel giúp detector
có đầu vào khác 0 để chọn offset xác định.

\subsection{LOD/LZD: quy ước, cấu trúc và vị trí sử dụng}

Trong thiết kế phần cứng Posit \cite{jaiswal2018posit,murillo2020flopoco}, bộ đếm
bit~0 dẫn đầu (LZD) và bộ đếm bit~1 dẫn đầu (LOD) thường được sử dụng song song
để xác định độ dài chuỗi \textit{regime} (lần lượt tìm bit~1 hoặc bit~0 đầu
tiên sau chuỗi lặp) cũng như chuẩn hóa kết quả cộng/trừ. Trong RTL của đề tài,
khối \sig{lod} được chuẩn hóa theo quy ước tìm \emph{vị trí bit 1 đầu tiên} từ MSB.
Khi \sig{vld=1}, vị trí $K$ tính từ 0 cũng chính là số bit 0 dẫn đầu
(LZC). Với đầu vào toàn 0, bộ dò trả $K=0$, \sig{vld=0}.

Parser chuyển cả hai bài toán đếm run về cùng một phép dò:
XOR hai bit kề tạo bit 1 tại chỗ đổi cực tính đầu tiên, dù
run bắt đầu bằng 0 hay 1. Khi đó độ dài run là $K+1$.
Nếu không có chuyển tiếp, sentinel làm $K=30$ và
\sig{vld=1}. Cách này dùng một đường dò chung cho hai cực tính.

Hình~\ref{fig:lod-regime-tree} khai triển cây gộp radix-4 trong
Parser A/B/C và Parser của quire. Hình~\ref{fig:lod-sac-adder}
cho thấy hai cách cài đặt khác: SAC dùng cây nhị phân bốn cấp
cho fraction 12 bit; S4 và Parser tích NF dùng \sig{casez} 31 bit.
Tên file \sig{lod\_lzd\_core\_fixed2} và \sig{fixed4}
không biểu thị radix. Đường READ của quire có phép dò riêng trên
word 36 bit (Hình~\ref{fig:quire-context}).

\rtlfig{LOD tìm chuyển tiếp regime của Parser A/B/C và quire: ba cấp gộp 4:1. Nhánh đệm hằng được vẽ để thể hiện cấu trúc khai triển của RTL.}{fig:lod-regime-tree}{
\node[pin] (in) at (0,0) {$d[30:0]=\{m[29:0]\oplus m[30:1],1'b1\}$};
\node[unit,text width=115mm] (pad) at (0,-1.0) {\hw{Căn trái vector tìm kiếm}\\
  \sig{padded[63:0]} $=\{d[30:0],33'b0\}$};
\draw[wire] (in)--(pad);
\foreach \n/\x/\r/\s in {3/-5.4/63:48/16 bit dữ liệu,2/-1.8/47:32/15 bit + 0,1/1.8/31:16/16 bit 0,0/5.4/15:0/16 bit 0}{
  \node[pad,text width=27mm] (g\n) at (\x,-2.68) {\sig{[\r]}\\ \s};
  \draw[wire] (pad.south -| g\n)--(g\n);
  \node[unit,text width=27mm] (a\n) at (\x,-4.24) {\hw{Cấp 1}\\$4\times$ gộp 4:1};
  \draw[wire] (g\n)--(a\n);
  \node[unit,text width=27mm] (b\n) at (\x,-5.93) {\hw{Cấp 2}\\gộp 4:1};
  \draw[wire] (a\n)--node[tag,right]{$4\times(v,p_2)$}(b\n);
}
\node[unit,text width=133mm] (root) at (0,-7.64) {\hw{Cấp 3: chọn nhóm hợp lệ cao nhất}\\
   \sig{pos[5:0]} $=\{g_3[1:0],p_4[3:0]\}$};
\foreach \n in {0,1,2,3}{\draw[wire] (b\n)--node[tag,right]{$(v,p_4)$}(root.north -| b\n);}
\node[pin] (out) at (0,-8.91) {\sig{vld} $=v_{\mathrm{root}}$ \qquad \sig{K[4:0]} $=$ \sig{vld ? pos[4:0] : 0}};
\draw[wire] (root)--(out);
\node[explain,text width=133mm] at (0,-10.25) {\hw{Một nút gộp}\\
 $v=\bigvee_{i=0}^{3}v_i$. Ưu tiên nhóm $3,2,1,0$; mã vị trí nhóm lần lượt $00,01,10,11$.\\
 Ghép mã nhóm ở phía cao và vị trí con được chọn ở phía thấp.};
\node[pad,text width=133mm] at (0,-12.1) {\hw{Ví dụ từ MSB: cùng detector cho hai dạng terminator}\\
 \begin{tabular}{@{}lll@{}}
   Đếm bit 1:  \sig{11101} & $\to$ tìm bit 0 đầu tiên & \sig{XOR = 0011}, $K=2$ \\
   Đếm bit 0:  \sig{00010} & $\to$ tìm bit 1 đầu tiên & \sig{XOR = 0011}, $K=2$
 \end{tabular}\\[2pt]
 Không có chuyển tiếp: sentinel cho $K=30$, \sig{vld=1}.};
}


\rtlfig{Hai bộ dò còn lại trong datapath NF/F: cây nhị phân 12 bit cho SAC và mã hóa ưu tiên 31 bit cho adder và Parser tích NF.}{fig:lod-sac-adder}{
\node[head] at (-3.78,0) {\sig{lod\_lzd\_core\_fixed4}\\SAC: dò trên mask fraction};
\node[head] at (3.78,0) {\sig{lod\_lzd\_core\_fixed2}\\Adder S4 và giải mã tích NF};
\node[pin] (fx) at (-3.78,-.97) {\sig{fx[11:0]}};
\node[pin] (in) at (3.78,-.97) {\sig{in[30:0]}};
\node[unit,text width=60mm] (pad) at (-3.78,-2.02) {\hw{16 lá bit}\\\sig{data\_pad = \{fx,4'b0\}}\\mỗi lá: $(v,p)=(\text{bit},0)$};
\draw[wire] (fx)--(pad);
\foreach \l/\y/\n/\step in {1/-3.49/8/1,2/-4.77/4/2,3/-6.05/2/4,4/-7.33/1/8}{
  \node[unit,text width=60mm] (b\l) at (-3.78,\y) {\hw{Cấp \l: \n\ cặp gộp} \quad \sig{STEP}=$\step$};
}
\draw[wire] (pad)--(b1);
\draw[wire] (b1)--(b2);
\draw[wire] (b2)--(b3);
\draw[wire] (b3)--(b4);
\node[explain,text width=60mm] (merge) at (-3.78,-8.96) {\hw{Quy tắc tại một cặp}\\
 $v=v_H\lor v_L$;\quad $p=v_H?p_H:(p_L\mathbin{|}\mathrm{STEP})$};
\draw[wire] (b4)--node[tag,right]{$(v,p)$}(merge);
\node[pin] (out1) at (-3.78,-10.34) {\sig{vld}, \sig{K[3:0]} $=\sig{vld}?p:0$};
\draw[wire] (merge)--(out1);
\node[unit,text width=60mm] (v) at (3.78,-2.02) {\hw{OR 31 bit}\\\sig{vld = |in}};
\node[unit,text width=60mm] (pz) at (3.78,-4.72) {\sig{casez}:~\hw{mẫu từ MSB}\\[1pt]
 \sig{1???...} $\to K=0$\\
 \sig{01??...} $\to K=1$\\
 \sig{001?...} $\to K=2$\\
 \sig{...}\\
 \sig{00...001} $\to K=30$\\[1pt]
 mặc định $K=0$};
\draw[wire] (in.south)--(v.north);
\draw[wire] (in.east)--++(2.9,0)|-(pz.east);
\node[unit,text width=60mm] (out2) at (3.78,-7.25) {\hw{Hai đầu ra tổ hợp}\\\sig{vld}, \sig{K[4:0]}};
\draw[wire] (pz)--node[tag,right]{$K$}(out2);
\draw[wire] (v.west)--++(-.37,0)|-(out2.west) node[pos=.46,tag,right]{\sig{vld}};
\node[pad,text width=60mm] at (3.78,-8.96) {\hw{Cùng mã hóa, hai nguồn vào}\\
 S4: \sig{magnitude[30:0]}\\
 NF: vector XOR kề + sentinel};
\node[tag] at (3.78,-10.34) {Ví dụ: \sig{in=0} $\Rightarrow$ \sig{vld=0, K=0}};
}



\subsection{P2 căn payload và tạo hai phần fraction}

P2 vừa dịch seed sang trái, vừa tạo regime signed từ \texttt{cnt}
và polarity. Khi regime bắt đầu bằng 1, $k=\texttt{cnt}$; khi bắt
đầu bằng 0, $k=-\texttt{cnt}-1$. Ba bit trên của payload dịch là
exponent. Phần còn lại được tách thành $f_H[11:0]$ và $f_L[13:0]$.
P2 ghi cả hai phần và \texttt{tail\_cut} bằng OR của $f_L$.
Ở A/B, phần thấp chỉ đóng góp thông tin cắt. Ở C, hai phần được
ghép lại để bộ cộng nhận đủ độ chính xác khả dụng.

Ví dụ FT-05, \texttt{1c900000} giải mã $sf_A=-10$ và
$F_A=\texttt{0x480}$. \texttt{3c820000} giải mã $sf_B=-1$
và $F_B=\texttt{0x208}$. Hai mã cùng dấu dương và phần thấp đã
bỏ bằng 0. Bộ chọn OPS thấy cả hai fraction có hai bit 1, nên
chọn A làm X và B làm Y. Cách kiểm các số này là giải mã raw
độc lập, sau đó so dấu, tổng scale và magnitude ở trace core.

\subsection{Điều kiện advance qua P1/P2}

Nếu P2 rỗng hoặc downstream nhận dữ liệu, P2 có thể chốt payload
mới từ P1. P1 có thể nhận thêm khi P1 rỗng hoặc khi P2 tiến.
Thanh ghi valid tiến cùng với dữ liệu. Đối với top hiện hành,
ba Parser retire cùng sự kiện \texttt{pair\_retire}; vì vậy
dữ liệu C không được ghép với tích từ request khác. Zero và NaR
đi theo metadata đến wrapper. Fraction hoặc scale của một mã
đặc biệt chỉ được sử dụng theo nhánh xử lý đã chọn.


\section{Bộ nhân FULL\_TERM}
\begin{reportpseudocode}{alg:full-term}{Nhân lặp dịch--cộng FULL\_TERM26/Q24 (Iterative Multiplication)}
  {Hai toán hạng giải mã $A,B$ (gồm $s,sf,f_H,z,\eta,\mathrm{cut}$), cờ $\mathrm{OPS}$ và giới hạn vòng $n$.}
  {Độ lớn chuẩn hóa $M_{\mathrm{norm}}[25:0]$ (Q24), dấu $s$, tỉ lệ $sf_P$, số vòng $t$ và các cờ cắt đuôi.}
  \begin{tabularx}{\linewidth}{@{}r >{\raggedright\arraybackslash}X >{\raggedright\arraybackslash}p{.33\linewidth}@{}}
    \algline{1}{$\mathrm{bypass}\gets z_A\lor z_B\lor \eta_A\lor \eta_B\lor \eta_C$}{Phát hiện Zero/NaR để bỏ qua vòng lặp.}
    \algline{2}{\textbf{if} $\mathrm{bypass}$ \textbf{then} $(X,Y)\gets (0,0)$ \textbf{else} $\mathrm{OPS\_Select}(A,B,\mathrm{OPS})$}{Chọn toán hạng có popcount nhỏ hơn làm $X$.}
    \algline{3}{$s\gets s_A\oplus s_B$, $b\gets sf_A+sf_B$, $\mathrm{in\_cut}\gets \mathrm{cut}_A\lor \mathrm{cut}_B$}{Tính dấu tích, tỉ lệ cơ sở và cờ cắt ngõ vào.}
    \algline{4}{$y_{\mathrm{base}}\gets \{1'b0, 1'b1, f_{H,Y}, 12'b0\}$}{Tạo hạng Q24 nền có bit ẩn $1$.}
    \algline{5}{$M\gets y_{\mathrm{base}}$, $fx\gets f_{H,X}$, $S\gets 0$, $t\gets 0$, $\mathrm{first}\gets 1$}{Khởi tạo thanh ghi tích lũy và mặt nạ SAC.}
    \algline{6}{\textbf{while} $fx\ne 0$ \textbf{and} $t<n$ \textbf{do}}{Dừng sớm khi hết bit 1 hoặc chạm ngưỡng $n$.}
    \algline{7}{\quad $c\gets \mathrm{LZC}(fx[11:0])$, $S\gets S+c+1$}{Tính bước nhảy và độ dịch $S$ từ điểm nhị phân.}
    \algline{8}{\quad $\mathrm{last\_one}\gets ((fx\mathbin{\&}(fx-1))=0)$}{Kiểm tra bit 1 hiện tại có phải bit cuối.}
    \algline{9}{\quad $M\gets (\mathrm{first}\ ?\ y_{\mathrm{base}}\ :\ M)+(y_{\mathrm{base}}\gg S)$}{Dịch phải hạng nền theo $S$ và cộng dồn.}
    \algline{10}{\quad $fx\gets \{fx[10:0], 1'b0\}\ll c$, $t\gets t+1$, $\mathrm{first}\gets 0$}{Loại bỏ bit 1 vừa xử lý bằng dịch trái tăng dần.}
    \algline{11}{\textbf{end while}}{}
    \algline{12}{\textbf{if} $M[25]=1$ \textbf{then} $sf_P\gets b+1$, $M_{\mathrm{norm}}\gets \{M[24:0],1'b0\}$}{Chuẩn hóa sau tích nếu tràn sang bit $25$.}
    \algline{13}{\textbf{else} $sf_P\gets b$, $M_{\mathrm{norm}}\gets \{M[23:0],2'b0\}$}{Ngược lại dịch trái 2 bit về dạng chuẩn 26 bit.}
    \algline{14}{$\mathrm{approx\_cut}\gets (fx\ne 0)$ \textbf{if} $\neg(z_P\lor \eta_P)$ \textbf{else} $0$}{Ghi nhận cờ cắt gần đúng nếu còn bit chưa cộng.}
    \algline{15}{\textbf{return} $M_{\mathrm{norm}}, s, sf_P, t, \mathrm{in\_cut}, \mathrm{approx\_cut}$}{Xuất tích chuẩn hóa sang tầng ghép MAC.}
  \end{tabularx}
\end{reportpseudocode}
\begin{samepage}
Wrapper xử lý riêng các trường hợp bypass Zero/NaR.
Ở RTL, context giữ $Y$, cấu hình và metadata. T1 phát vị trí dịch và token
first/last. T2 dịch để tạo term. T3 cộng vào accumulator, chọn $M_0$ làm
giá trị nền ở token đầu. Token init-only có term bằng 0. Tín hiệu
\texttt{core\_done} xuất hiện cùng kết quả commit của token cuối.
Ví dụ FT-05 với A=\texttt{1c900000}, B=\texttt{3c820000} (hex),
OPS=1 và n ít nhất 2, cho $M=24\,224\,768$, $b=-11$, $t=2$.
\par\end{samepage}

\subsection{OPS và context đầu bộ nhân}

OPS đếm bit 1 trên phần fraction 12 bit của A và B. Mỗi phép
đếm được chia thành hai nhóm 6 bit, rồi cộng hai số đếm cục bộ.
OPS1 chọn B làm X chỉ khi count B nhỏ hơn count A. Trường hợp
hòa giữ A làm X. Do X quyết định số token cần phát, OPS1 có thể
giảm số vòng cho cùng cặp A/B. OPS0 giúp kiểm vai cố định.
Y gồm bit ẩn và phần fraction 12 bit, nên bus Y có 13 bit.
Dấu tích là XOR, số mũ cơ sở là $sf_A+sf_B$ signed 11 bit.

\rtlfig{OPS, SAC và SBM trong bộ nhân FULL\_TERM26/Q24. T1/T2 chốt token; T3 chốt tổng. Hai đường hồi tiếp phục vụ cập nhật mask và cộng term.}{fig:multiplier-stages}{
\node[pin] (ab) at (-1.55,.46) {$f_A[11:0], f_B[11:0]$};
\node[pin] (cfg) at (5.4,.46) {\sig{ops\_q}, \sig{n\_q[3:0]}};
\node[unit,text width=74mm] (ops) at (-1.55,-1.06) {\hw{OPS}\\
 mỗi toán hạng: 2 đếm 6 bit $\to$ tổng 4 bit\\
 B làm X nếu OPS1 và $\mathrm{popcount}(B)<\mathrm{popcount}(A)$};
\draw[wire] (ab)--(ops);
\draw[enable] (cfg.south)|-node[tag,above]{OPS}(ops.east);
\node[ff,text width=64mm] (sacreg) at (-2.45,-3.02) {\hw{Context SAC tại launch}\\
 \sig{fx\_q[11:0]} $\leftarrow X$; scale/count $\leftarrow0$\\
 \sig{limit\_q[3:0]} $\leftarrow n$};
\node[ff,text width=47mm] (y) at (4.55,-3.02) {\hw{Context Y và metadata}\\\sig{y\_q[12:0]} $\leftarrow \{1,f_Y\}$\\dấu, $sf$, \sig{input\_cut}};
\draw[wire] (ops.south -| sacreg)--node[tag,right]{X:12}(sacreg);
\draw[wire] (ops.south)--++(0,-.31)-|(y.north) node[pos=.25,tag,above]{Y:13};
\node[unit,text width=64mm] (sac) at (-2.45,-5.18) {\hw{SAC + điều khiển phát token}\\
 LOD12 $\to (K,v)$; $S' = S+K+1$\\
 $F'=\{F[10:0],0\}\ll K$ khi $v=1$\\
 hết bit nếu $F\mathbin{\&}(F-1)=0$};
\draw[wire] (sacreg)--node[tag,right]{$F,S$, count, limit}(sac);
\draw[wire] (sac.west)--++(-.52,0)|-(sacreg.west);
\node[explain,text width=46mm] at (4.55,-5.2) {\hw{Dừng phát}\\init nếu $F=0$ hoặc $n=0$.\\
 last nếu init, hết bit,\\hoặc count+1 đạt $n$.};
\node[ff,text width=64mm] (t1) at (-2.45,-7.25) {\hw{T1} \quad scale 4, iteration 4\\valid, first, last, init, cut};
\draw[wire] (sac)--(t1);
\node[unit,text width=45mm] (base) at (4.55,-7.25) {\hw{Cơ số Q24}\\$B_0=\{0,\sig{y\_q},12'b0\}$};
\draw[wire] (y.east)--++(.33,0)|-(base.east);
\node[unit,text width=64mm] (sbm) at (-2.45,-8.98) {\hw{SBM: dịch phải 26 bit}\\
 các cấp 1/2/4/8; term $=B_0\gg S'$\\init-only chọn term $=0$};
\draw[wire] (t1)--(sbm);
\draw[wire] (base.south)--++(0,-.35)|-(sbm.east) node[pos=.63,tag,above]{26};
\node[ff,text width=64mm] (t2) at (-2.45,-10.78) {\hw{T2} \quad term 26, iteration 4\\valid, first, last, cut};
\draw[wire] (sbm)--(t2);
\node[pick,text width=32mm] (m) at (4.55,-11.04) {first ? $B_0$ : acc};
\draw[wire] (base)--(m);
\node[unit,text width=64mm] (sum) at (-2.45,-12.53) {\hw{Cộng mở rộng 27 bit}\\base + term; bit 26 kiểm carry};
\draw[wire] (t2)--node[tag,right]{term:26}(sum);
\draw[wire] (m.south)|-(sum.east) node[pos=.62,tag,above]{base:26};
\draw[enable] (t2.east)|-(m.west) node[midway,tag,above]{first};
\node[ff,text width=64mm] (t3) at (-2.45,-14.25) {\hw{T3: commit khi T2 valid}\\acc 26, iteration 4, cut, done};
\draw[wire] (sum)--(t3);
\draw[wire] (t3.east)--(7.45,-14.25) |- (m.east) node[pos=.61,tag,left]{acc:26};
\node[pin] (out) at (-2.45,-15.45) {\sig{acc[25:0]}, \sig{iterations\_done}, \sig{approx\_cut}, \sig{core\_done}};
\draw[wire] (t3)--(out);
}



\subsection{SAC, SBM và tích lũy không mất bit}

SAC dùng mask còn lại để tìm bit 1 có trọng số cao nhất, tạo
offset và mask cho bước tiếp theo. T1 chốt offset tích lũy,
token first/last, số vòng và cờ cắt. Token last được tạo khi
mask cạn hoặc khi đã dùng hết số vòng. Nếu không xử lý bit
fraction nào, token init-only chuyển tới T2 với term bằng 0.
Nhờ vậy đường dữ liệu commit vẫn chung với trường hợp có vòng.

T2 dịch phải cơ số \texttt{y\_base} theo offset của token.
Các vị trí bit thấp được đệm 0 từ trước, nên phép dịch từ
1 đến 12 giữ đủ số hạng trong lưới Q24. T3 chọn cơ số Y ở token
first hoặc accumulator ở các token sau, cộng term và lưu kết
quả. \texttt{core\_done} cùng có hiệu lực với kết quả của token
last đã commit. Số vòng fraction được cập nhật từ metadata của
token T2, không đếm slot init-only.

\begin{engineeringtable}{Trace số học bộ dữ liệu mẫu FT-05 trên MAC-F ở cấu hình OPS1/n12, $C=0$.}{tab:ft05-terms}
\begin{tabularx}{\linewidth}{@{}r l r Y@{}}
\toprule
Cạnh & Mốc & $M$ quan sát & Diễn giải \\
\midrule
0 & Accept & -- & Chốt A/B và cấu hình. \\
2 & \texttt{pair\_retire} & -- & Context core và C được chốt. \\
5 & T2 valid, vòng 1 & -- & T3 sắp commit số hạng đầu. \\
6 & T2 valid, vòng 2 & 23.633.920 & Cơ số 18.907.136 cộng term dịch phải 2. \\
7 & \texttt{core\_done} & 24.224.768 & Cộng thêm 590.848 từ dịch phải 5, $b=-11$, $t=2$. \\
14 & \texttt{out\_valid} & 24.224.768 & $D=\texttt{1ae34800}$, flags \texttt{00}. \\
\bottomrule
\end{tabularx}
\end{engineeringtable}

Ở n=0 với cùng A/B, init-only cho $M=18\,907\,136$, $t=0$ và
cờ còn số hạng được đặt. Mã đo được là \texttt{1a410000} với
flags \texttt{03}. Ở n12, X cạn sau hai vòng. Hai trường hợp
cho thấy thời gian xử lý phụ thuộc số bit 1 được chọn và không
đồng nhất với giá trị giới hạn cấu hình. Do bộ nhân luôn chạy ít nhất một chu kỳ ở pha lặp \texttt{T2/T3} sau \texttt{T1}, số chu kỳ thực thi của khối nhân lặp bằng $\max(1,t)$ chu kỳ cộng pha khởi tạo cố định, quyết định trực tiếp thành phần biến thiên trong độ trễ từng phép MAC.

\subsection{Adapter ở biên trước bộ cộng}

Nếu $M[25]=1$, significand tích nằm từ 2 đến dưới 4, nên scale
chuẩn hóa là $b+1$ và fraction lấy từ \texttt{\{M[24:0],1'b0\}}.
Nếu bit 25 bằng 0, scale giữ $b$ và fraction là
\texttt{\{M[23:0],2'b0\}}. Wrapper tính trước $b+1$ khi chốt
context, để adapter chọn giữa hai scale thay vì thêm carry-chain
vào đường từ core đến S1. Cả hai cách ánh xạ giữ đủ bit có giá
trị trong tích Q24. NF sử dụng biểu diễn đã chuẩn hóa để qua
Packer tích và Parser tích trước khi cộng, còn F dùng trực tiếp.


\section{Cộng, chuẩn hóa và đóng gói}
Adapter của MAC-F tăng số mũ khi bit 25 của accumulator bằng 1 và
căn lại fraction 26 bit theo bit dẫn đầu. Cách ánh xạ này giữ đủ dữ liệu
tích. Bộ cộng đưa hai toán hạng về cùng thang, cộng hoặc trừ theo dấu và
chuẩn hóa tổng.

\begin{table}[!htbp]
  \centering
  \caption{Năm tầng của bộ cộng dùng chung.}
  \label{tab:adder-stages}
  \begin{tabular}{lp{0.74\textwidth}}
    \hline
    Tầng & Nhiệm vụ \\
    \hline
    S1 & Chọn toán hạng có độ lớn lớn hơn và tính chênh lệch số mũ. \\
    S2 & Dịch phải để căn thang. Gộp phần đuôi bị dịch thành jam/sticky. \\
    S3 & Cộng hoặc trừ magnitude 31 bit, tính trước scale khi có carry. \\
    S4 & Mã hóa ưu tiên bit 1 dẫn đầu trên 31 bit và tạo seed chuẩn hóa. \\
    S5 & Dịch chuẩn hóa, cập nhật scale và trả metadata ngoại lệ. \\
    \hline
  \end{tabular}
\end{table}

Packer P1 dựng chuỗi regime, exponent và fraction cùng guard, round,
sticky. P2 quyết định RNE, áp dụng kẹp và tạo mã dấu theo hợp đồng.
Giữ dữ liệu tích tới bộ cộng giúp MAC-F tránh lượt RNE tích trung gian.
Các trường hợp gần triệt tiêu vẫn nhạy với sai số đã phát sinh khi cắt
fraction đầu vào và với bước lượng tử hóa cuối.

\subsection{Cấu trúc đường ống 5 tầng co giãn của \texttt{posit\_add}}

Bộ cộng dấu phẩy động Posit (\texttt{posit\_add}) được tổ chức thành đường ống
5 tầng co giãn \texttt{S1..S5} nhằm cân bằng độ sâu logic ở mọi tầng và đáp ứng
mục tiêu tần số từ $100\text{ MHz}$ trở lên trên Artix-7: mỗi tầng chỉ chứa tối
đa một phép toán số học hoặc một bộ dịch lớn, giúp năm tầng \texttt{S1..S5}
đạt độ trễ tổ hợp sau tổng hợp nằm trong khoảng $2{,}958\text{--}4{,}851\text{ ns}$
(Bảng~\ref{tab:f2f-stages}) và đóng định thời định tuyến trên MAC-F ở
$F_{\max,\mathrm{R2R}}=154{,}42\text{ MHz}$ (Bảng~\ref{tab:fpga-synth-vs-route}).
Toàn bộ mạch tuần tự và tín hiệu bắt tay \texttt{valid/ready} được tập trung tại
mô-đun đỉnh \texttt{posit\_add.sv}, bao bọc năm khối tổ hợp thuần túy từ S1 đến S5
(Hình~\ref{fig:adder-stages}).

Ở cơ chế co giãn, tín hiệu \texttt{ready} truyền ngược từ S5 về S1: một slot rỗng
có thể nhận ngay, còn slot đang giữ dữ liệu hợp lệ chỉ tiến lên khi tầng kế tiếp
sẵn sàng; nếu bị chặn (\texttt{out\_ready=0}), toàn bộ dữ liệu số học và
metadata (dấu, cờ \texttt{Zero/NaR}, phần dư và \texttt{flags}) được giữ đồng bộ
với nhau vì chỉ lệch một trường cũng làm sai lệch phép đóng gói cuối.

\rtlfig{Datapath mô-đun \sig{posit\_add}. Cột trái là xử lý tổ hợp; cột phải liệt kê các trường chính được chốt sau mỗi tầng.}{fig:adder-stages}{
\node[pin] (in) at (-2.3,.28) {Hai số hạng: frac 26, $sf$ 11, $s,Z,N$ mỗi số};
\node[head] at (4.65,0) {Ngân hàng thanh ghi};
\node[unit,text width=73mm] (s1) at (-2.3,-1.42) {\hw{S1: so và đổi vai}\\
 ưu tiên số khác 0; so $sf$, rồi so frac\\
 hai hiệu scale signed 12 bit $\to |\Delta sf|$; kẹp nếu $\Delta>28$};
\node[ff,text width=50mm] (r1) at (4.65,-1.42) {\hw{S1}\\frac X/Y: 26+26; shift: 5\\kẹp, subtract, $sf_{\mathrm{base}}$, dấu/cờ};
\draw[wire] (in)--(s1);
\draw[wire] (s1)--(r1);
\node[unit,text width=73mm] (s2) at (-2.3,-3.96) {\hw{S2: căn Y vào cửa sổ 31 bit}\\
 $\{0,1,f_Y,3'b0\}\gg\mathrm{shift}$\\
 mặt nạ đuôi $\to$ jam ở \sig{Y[0]} và \sig{alignment\_tail}};
\node[ff,text width=50mm] (r2) at (4.65,-3.96) {\hw{S2}\\frac X: 26; Y: 31; tail: 1\\subtract, $sf_{\mathrm{base}}$, dấu/cờ};
\draw[wire] (r1.south)--++(0,-.54)-|(s2.north);
\draw[wire] (s2)--(r2);
\node[unit,text width=73mm] (s3) at (-2.3,-6.51) {\hw{S3: cộng/trừ độ lớn}\\
 $X=\{0,1,f_X,3'b0\}$; $M=X\pm Y$\\
 song song tính \sig{sf\_inc} $=sf_{\mathrm{base}}+1$};
\node[ff,text width=50mm] (r3) at (4.65,-6.51) {\hw{S3}\\magnitude: 31; tail: 1\\\sig{sf\_inc}: 11, dấu/cờ};
\draw[wire] (r2.south)--++(0,-.54)-|(s3.north);
\draw[wire] (s3)--(r3);
\node[unit,text width=73mm] (s4) at (-2.3,-9.06) {\hw{S4: LOD ưu tiên và seed}\\
 $M[30:0]\to$ \sig{leading\_count[4:0]}, \sig{nonzero}\\
 song song ghép seed 30 bit và \sig{sticky\_raw}};
\node[ff,text width=50mm] (r4) at (4.65,-9.06) {\hw{S4}\\seed: 30; count: 5; nonzero\\sticky, \sig{sf\_inc}, dấu/cờ};
\draw[wire] (r3.south)--++(0,-.54)-|(s4.north);
\draw[wire] (s4)--(r4);
\node[unit,text width=73mm] (s5) at (-2.3,-11.61) {\hw{S5: chuẩn hóa và phân loại}\\
 dịch trái seed theo count; lấy 29 bit cao\\
 $sf=\sig{sf\_inc}-\mathrm{count}$; xác định Zero/NaR};
\node[ff,text width=50mm] (r5) at (4.65,-11.61) {\hw{S5}\\\sig{frac\_hi[28:0]}, $sf$, sticky\\dấu, Zero/NaR và flags};
\draw[wire] (r4.south)--++(0,-.54)-|(s5.north);
\draw[wire] (s5)--(r5);
\node[pin] (o) at (4.65,-13.08) {Tới Packer: \sig{frac = \{frac\_hi,24'b0\}}};
\draw[wire] (r5)--(o);
}



Toàn bộ luồng tính toán năm tầng của \texttt{posit\_add} được tóm tắt
trong Algorithm~\ref{alg:posit-add}, đối chiếu trực tiếp với các khối tổ
hợp từ S1 đến S5 trong RTL.

\begin{reportpseudocode}{alg:posit-add}{Cộng Posit(32,3) đường ống 5 tầng co giãn (Five-Stage Elastic Addition)}
  {Hai số hạng chuẩn hóa $A,B$ (với $s, sf[10:0], f[25:0], z, \eta$) và cờ $\mathrm{flags_{in}}[4:0]$.}
  {Kết quả chuẩn hóa $(s_{\mathrm{out}}, sf[10:0], f_{\mathrm{out}}[52:0], \mathrm{sticky}, z_{\mathrm{out}}, \eta, \mathrm{flags_{out}}[4:0])$.}
  \begin{tabularx}{\linewidth}{@{}r >{\raggedright\arraybackslash}X >{\raggedright\arraybackslash}p{.33\linewidth}@{}}
    \algline{1}{\textbf{S1} \quad $d_{AB}\gets sf_A-sf_B$, $d_{BA}\gets sf_B-sf_A$}{Tính hai hiệu scale mở rộng dấu 12 bit song song.}
    \algline{2}{$\mathrm{choose\_a}\gets \neg z_A\land \bigl(z_B\lor (d_{AB}\ge 0 \land (sf_A\ne sf_B \lor f_A\ge f_B))\bigr)$}{Chọn toán hạng có độ lớn lớn hơn (ưu tiên khác 0).}
    \algline{3}{$(X,Y)\gets (A,B)$ \textbf{if} $\mathrm{choose\_a}$ \textbf{else} $(B,A)$, $\delta \gets \min(|\Delta sf|,29)$}{Hoán đổi toán hạng; xác định độ dịch $\delta$ và cờ kẹp $|\Delta sf|>28$.}
    \algline{4}{$\mathrm{sub}\gets \neg(z_A\lor z_B)\land (s_A\ne s_B)$, $\eta\gets \eta_A\lor\eta_B\lor\mathrm{flags_{in}}[4]$}{Xác định phép trừ hiệu dụng và lan truyền cờ NaR.}
    \algline{5}{\textbf{S2} \quad $w_Y\gets \{1'b0, 1'b1, f_Y[25:0], 3'b0\}$ \textbf{if} $\neg(z_A\lor z_B)$ \textbf{else} $0$}{Mở rộng toán hạng nhỏ thành từ 31 bit có bit ẩn $1$.}
    \algline{6}{$y_{\mathrm{align}}[30:1]\gets (w_Y\gg \delta)[30:1]$, $\mathrm{tail}\gets \mathrm{clamped}\lor \mathrm{raw\_tail}$}{Dịch phải 31 bit và gom mặt nạ nhiệt kế cho bit rơi.}
    \algline{7}{$y_{\mathrm{align}}[0]\gets ((w_Y\gg \delta)[0]\lor \mathrm{tail})$ \textbf{if} $\neg(z_A\lor z_B)$ \textbf{else} $0$}{Chèn bit jam tại LSB khi dịch rơi bit hoặc kẹp $\delta$.}
    \algline{8}{\textbf{S3} \quad $w_X\gets \{1'b0, 1'b1, f_X[25:0], 3'b0\}$ \textbf{if} $\neg(z_A\land z_B)$ \textbf{else} $0$}{Đưa toán hạng lớn về dạng 31 bit để cộng/trừ.}
    \algline{9}{$M\gets w_X-y_{\mathrm{align}}$ \textbf{if} $\mathrm{sub}$ \textbf{else} $w_X+y_{\mathrm{align}}$, $sf_{\mathrm{inc}}\gets sf_X+1$}{Cộng/trừ 31 bit và cộng bù trước $1$ cho scale.}
    \algline{10}{\textbf{S4} \quad $(K,\mathrm{nz})\gets \mathrm{LOD}(M[30:0])$, $\mathrm{lost}\gets M[30]\land M[0]$}{Dò bit 1 dẫn đầu thống nhất 31 bit; kiểm tra carry bị mất.}
    \algline{11}{$\mathrm{seed}[29:0]\gets \{M[29:2], M[1]\lor\mathrm{lost}, M[0]\}$, $\mathrm{sticky}_{\mathrm{raw}}\gets \mathrm{tail}\lor\mathrm{lost}$}{Tạo hạt dịch 30 bit và cập nhật sticky thô.}
    \algline{12}{\textbf{S5} \quad $f_{\mathrm{hi}}[28:0]\gets (\mathrm{seed}\ll K)[29:1]$, $sf\gets sf_{\mathrm{inc}}-K$}{Dịch trái 30 bit chuẩn hóa và cập nhật scale đồng nhất.}
    \algline{13}{$z_{\mathrm{out}}\gets \neg\eta\land ((z_A\land z_B)\lor \neg\mathrm{nz})$, $f_{\mathrm{out}}\gets \{f_{\mathrm{hi}}, 24'b0\}$}{Phân loại Zero/NaR và nối 24 bit 0 tạo ngõ ra 53 bit.}
    \algline{14}{\textbf{return} $s_{\mathrm{out}}, sf, f_{\mathrm{out}}, \mathrm{sticky}_{\mathrm{raw}}, z_{\mathrm{out}}, \eta, \mathrm{flags}_{\mathrm{out}}$}{Chốt thanh ghi S5 kèm cập nhật bit \texttt{inexact}.}
  \end{tabularx}
\end{reportpseudocode}

\subsection{Tầng S1: Chuẩn bị, so sánh độ lớn và hoán đổi toán hạng}

Ở đầu vào MAC-F, $X$ và $Y$ của bộ cộng là hai số hạng được sắp lại theo độ lớn
($|X|\ge|Y|$); ký hiệu này dùng trong phạm vi bộ cộng và khác với vai $X/Y$ của
khối OPS ở bộ nhân. Mỗi số hạng đưa vào S1 gồm fraction 26 bit, scale signed 11 bit,
dấu và cờ đặc biệt. Tầng S1 (\texttt{posit\_add\_prepare\_comb}, Hình~\ref{fig:adder-s1})
tính song song hai hiệu scale mở rộng dấu 12 bit (\texttt{diff\_ab}, \texttt{diff\_ba}),
so sánh ngang bằng scale (\texttt{sf\_equal}) và so sánh fraction (\texttt{frac\_a\_ge\_b})
để tạo tín hiệu điều khiển \texttt{choose\_a\_cond} ưu tiên toán hạng khác 0.

Cờ kẹp \texttt{clamp\_ab/\!ba} kiểm tra đồng thời nhóm bit cao \texttt{diff[10:5]}
khác 0 hoặc \texttt{diff[4:0] > 28} để nhận biết $|\Delta sf|>28$ mà không bị
cuộn modulo 5 bit khi chênh lệch scale lớn. Để giảm fanout và tránh nghẽn định
tuyến trên đường tráo toán hạng $62$ nhánh, RTL tách tín hiệu chọn thành bốn lưới
có thuộc tính \texttt{dont\_touch}: \texttt{choose\_a\_frac0..2} điều khiển riêng
ba lát fraction \texttt{[8:0]}, \texttt{[17:9]}, \texttt{[25:18]} và
\texttt{choose\_a\_meta} điều khiển \texttt{sign\_base}/\texttt{sf\_base}, bảo toàn
nguyên vẹn ranh giới thanh ghi và số chu kỳ đường ống.

\rtlfig{Datapath chi tiết tầng S1 (\sig{posit\_add\_prepare\_comb}): hai bộ trừ scale 12 bit song song, phát hiện kẹp $|\Delta sf|>28$ và bốn nhánh \sig{choose\_a\_*} cân bằng tải.}{fig:adder-s1}{
\node[pin] (ina) at (-5.05,.82) {\sig{a\_sf},\,\sig{b\_sf[10:0]}};
\node[pin] (inb) at (.20,.82) {\sig{a\_frac},\,\sig{b\_frac[25:0]}};
\node[pin] (inc) at (5.35,.82) {dấu, cờ \sig{Z/N}, \sig{flags\_in[4:0]}};

\node[unit,text width=40mm] (sub) at (-5.05,-.92) {\hw{Trừ scale 12 bit song song}\\
 $\sig{diff\_ab}=\{a_{10},a_{sf}\}-\{b_{10},b_{sf}\}$\\
 $\sig{diff\_ba}=\{b_{10},b_{sf}\}-\{a_{10},a_{sf}\}$\\
 $\sig{sf\_a\_lt\_b}=\sig{diff\_ab}[11]$};
\node[unit,text width=45mm] (cmp) at (.20,-.92) {\hw{So sánh độ lớn toán hạng}\\
 $\sig{sf\_equal}=(a_{sf}==b_{sf})$\\
 $\sig{frac\_a\_ge\_b}=(a_{\mathrm{frac}}\ge b_{\mathrm{frac}})$};
\node[unit,text width=38mm] (exc) at (5.35,-.92) {\hw{Cờ ngoại lệ \& phép tính}\\
 $\sig{both\_nz}=\neg a_Z\land\neg b_Z$\\
 $\sig{sub}=\sig{both\_nz}\land(a_s\ne b_s)$\\
 \sig{both\_zero}, \sig{is\_nar}, \sig{flags\_lo}};

\draw[wire] (ina)--(sub);
\draw[wire] (inb)--(cmp);
\draw[wire] (inc)--(exc);

\node[unit,text width=40mm] (clp) at (-5.05,-3.92) {\hw{Phát hiện kẹp \& dịch}\\
 $\sig{clamp\_ab/\!ba}: |\Delta sf|>28$\\
 chọn theo \sig{sf\_a\_lt\_b}:\\
 \sig{shift\_amount[4:0]}, \sig{clamped}};
\node[pick,text width=47mm] (sel) at (.20,-3.92) {\hw{Chọn $X{\ge}Y$ \& tách 4 tải}\\
 $\sig{choose\_a\_cond}=\neg a_Z\land(b_Z\lor\dots)$\\
 4 lưới \sig{dont\_touch}:\\
 \sig{choose\_a\_frac0..2}, \sig{meta}};

\draw[wire] (sub)--(clp);
\draw[wire] (cmp)--(sel);

\coordinate (gx) at (-1.50,-2.65);
\node[tapdot] at (-5.05,-2.65) {};
\draw[meta] (-5.05,-2.65) -- (gx) -- (sel.north -| gx);
\node[tag] at (-3.15,-2.65) {\sig{sf\_a\_lt\_b}};

\coordinate (gz) at (1.90,-2.65);
\node[tapdot] at (5.35,-2.65) {};
\draw[meta] (5.35,-2.65) -- (gz) -- (sel.north -| gz);
\node[tag] at (3.65,-2.65) {$a_Z, b_Z$};

\node[unit,text width=51mm] (muxxy) at (.20,-6.22) {\hw{MUX tráo 3 lát \& metadata}\\
 $[8:0]\gets\sig{frac0}$, $[17:9]\gets\sig{frac1}$, $[25:18]\gets\sig{frac2}$\\
 xuất \sig{x\_frac}, \sig{y\_frac}, \sig{sign\_base}, \sig{sf\_base}};
\draw[meta] (sel)--(muxxy);
\node[tag,anchor=west] at (.38,-5.10) {4 nhánh \sig{dont\_touch}};

\node[ff,text width=140mm] (r1) at (.15,-7.95) {\hw{S1}\enspace
 \sig{s1\_x\_frac\_q}, \sig{s1\_y\_frac\_q[25:0]} \enspace|\enspace
 \sig{s1\_shift\_q[4:0]}, \sig{s1\_clamped\_q}\\
 \sig{s1\_sub\_q}, \sig{s1\_sign\_q}, \sig{s1\_sf\_q[10:0]} \enspace|\enspace
 \sig{s1\_both\_nz\_q}, \sig{s1\_both\_zero\_q}, \sig{s1\_nar\_q}, \sig{s1\_flags\_q[3:0]}};

\draw[wire] (clp.south)--(r1.north -| clp);
\draw[wire] (muxxy.south)--(r1.north -| muxxy);
\draw[wire] (exc.south)--(r1.north -| exc);

\node[pin] (o) at (.15,-9.12) {Tới S2 (\sig{posit\_add\_align\_comb}): căn cửa sổ 31 bit, tạo \sig{jam} và \sig{alignment\_tail}};
\draw[wire] (r1)--(o);
}


\subsection{Tầng S2: Dịch phải căn chỉnh 31 bit và mặt nạ nhiệt tạo jam}

Tầng S2 (\texttt{posit\_add\_align\_comb}, Hình~\ref{fig:adder-s2}) ghép số hạng
nhỏ thành từ 31 bit $\texttt{y\_seed} = \{1'\mathrm{b0}, 1'\mathrm{b1}, \texttt{y\_frac}, 3'\mathrm{b000}\}$
và dịch phải qua bộ dịch \texttt{dyn\_right\_shifter\_fixed4} theo
\texttt{shift\_amount[4:0]}. Song song với bộ dịch, mạch tổ hợp gồm 25 bộ so sánh
hằng dựng mặt nạ nhiệt $\texttt{discard\_mask}[j] = (j + 3 < \texttt{shift\_amount})$
để thu gọn toàn bộ các bit bị dịch rơi xuống dưới bit~0 thành \texttt{raw\_tail}.

Khi \texttt{both\_nonzero=1} và không bị kẹp, 30 bit cao lấy phần dịch
\texttt{shifted\_full[30:1]}; ngược lại (\texttt{clamped=1} hoặc có một số hạng
bằng 0), 30 bit cao được xóa về 0. Bit \texttt{y[0]} được hợp nhất cùng phần dư
$\texttt{both\_nonzero} \land (\texttt{clamped} \lor \texttt{shifted\_full}[0] \lor \texttt{raw\_tail})$,
còn cờ \texttt{alignment\_tail} lưu riêng trạng thái rơi bit của phép căn chỉnh để
truyền tiếp qua S3 tới S4. Nhờ kiểm soát qua \texttt{both\_nonzero}, nhánh kẹp
không bao giờ tự phát sinh bit jam từ toán hạng 0.

\rtlfig{Datapath chi tiết tầng S2 (\sig{posit\_add\_align\_comb}): dịch phải căn chỉnh 31 bit, mặt nạ nhiệt 25 bit tính song song và hợp nhất \sig{jam}/\sig{alignment\_tail}.}{fig:adder-s2}{
\node[pin] (in1) at (-4.9,.86) {\sig{s1\_y\_frac\_q[25:0]}, \sig{s1\_shift\_q[4:0]}};
\node[pin] (in2) at (.2,.86) {\sig{s1\_clamped\_q}, \sig{s1\_both\_nz\_q}};
\node[pin] (in3) at (5.25,.86) {\sig{s1\_x\_frac\_q[25:0]}, \sig{s1\_sub\_q},\\ \sig{s1\_sf\_q[10:0]}, dấu, cờ};

\node[unit,text width=43mm] (sh) at (-4.9,-1.38) {\hw{Dịch phải 31 bit}\\
 \sig{dyn\_right\_shifter\_fixed4}\\
 $\sig{y\_seed}=\{0,1,f_Y,3'\mathrm{b0}\}$\\
 $\sig{shifted\_full}=\sig{y\_seed}\gg\mathrm{shift}$};
\node[unit,text width=44mm] (msk) at (.2,-1.38) {\hw{Mặt nạ nhiệt 25 bit song song}\\
 $\sig{discard\_mask}[j]=(j+3<\mathrm{shift})$\\
 $\sig{raw\_tail}=\bigvee(f_Y[24:0]\land\sig{mask})$};
\node[unit,text width=36mm] (pas) at (5.25,-1.38) {\hw{Giữ nguyên toán hạng X}\\
 \sig{x\_frac[25:0]}, \sig{subtract},\\
 \sig{sf\_base[10:0]}, \sig{sign\_base},\\
 \sig{both\_zero}, \sig{is\_nar}, \sig{flags\_lo}};

\draw[wire] (in1)--(sh);
\draw[wire] (in2)--(msk);
\draw[wire] (in3)--(pas);
\node[tapdot] at (-4.9,.18) {};
\draw[wire] (-4.9,.18) -- (-1.35,.18) -- (msk.north -| -1.35,0);
\node[tag] at (-2.9,.18) {\sig{f\_Y, shift}};

\node[unit,text width=43mm] (yhi) at (-4.9,-3.85) {\hw{Chốt 30 bit cao của Y}\\
 $\sig{y[30:1]}=(\sig{bnz}\land\neg\sig{clp})$\\
 $\text{ ? }\sig{shifted\_full[30:1]} : 30'\mathrm{d0}$};
\node[unit,text width=45mm] (jam) at (.2,-3.85) {\hw{Tạo jam Y[0] và tail}\\
 $\sig{y[0]}=\sig{bnz}\land(\sig{clp}\lor\sig{sh[0]}\lor\sig{raw\_tail})$\\
 $\sig{tail}=\sig{bnz}\land(\sig{clp}\lor\sig{raw\_tail})$};

\draw[wire] (sh)--(yhi);
\draw[wire] (msk)--(jam);
\coordinate (sh0p) at (-1.45,-2.66);
\node[tapdot] at (-4.9,-2.66) {};
\draw[wire] (-4.9,-2.66) -- (sh0p) -- (jam.north -| sh0p);
\node[tag] at (-3.05,-2.66) {\sig{sh[0]}};

\node[ff,text width=140mm] (r2) at (.15,-5.68) {\hw{S2}\enspace
 \sig{s2\_x\_frac\_q[25:0]} \enspace|\enspace
 \sig{s2\_y\_q[30:0]} \enspace|\enspace
 \sig{s2\_tail\_q}\\
 \sig{s2\_sub\_q}, \sig{s2\_sign\_q}, \sig{s2\_sf\_q[10:0]} \enspace|\enspace
 \sig{s2\_both\_zero\_q}, \sig{s2\_nar\_q}, \sig{s2\_flags\_q[3:0]}};

\draw[wire] (yhi.south)--(r2.north -| yhi);
\draw[wire] (jam.south)--(r2.north -| jam);
\draw[wire] (pas.south)--(r2.north -| pas);

\node[pin] (o) at (.15,-6.82) {Tới S3 (\sig{posit\_add\_arith\_comb}): cộng/trừ 31 bit \sig{X +/- Y} và tính song song \sig{sf\_inc}};
\draw[wire] (r2)--(o);
}


\subsection{Tầng S3 và S4: Cộng/trừ độ lớn, tiền tăng scale và dò bit dẫn đầu}

Hình~\ref{fig:adder-tail} minh họa đầy đủ đường dữ liệu của hai tầng S3
(\texttt{posit\_add\_arith\_comb}) và S4 (\texttt{posit\_add\_lod\_comb}):
\begin{itemize}
  \item \textbf{Tầng S3:} Ghép $X = \{1'\mathrm{b0}, 1'\mathrm{b1}, \texttt{x\_frac}, 3'\mathrm{b000}\}$
  rồi tính \texttt{magnitude[30:0] = subtract ? (X - Y) : (X + Y)}. Vì S1 đã bảo đảm
  $|X|\ge|Y|$ trên mọi cặp hữu hạn khác 0, hiệu $X-Y$ luôn không âm nên không cần
  mạch lấy trị tuyệt đối ở đầu ra S3. Song song với phép cộng/trừ, S3 tính trước
  $\texttt{sf\_inc} = \texttt{sf\_base} + 11'\mathrm{sd1}$ để chuẩn bị cho bước chuẩn hóa ở S5.
  \item \textbf{Tầng S4:} Bộ dò ưu tiên 31 bit \texttt{lod\_lzd\_core\_fixed2} quét từ
  MSB của \texttt{s3\_mag\_q[30:0]} để xác định \texttt{leading\_count[4:0]} và cờ
  \texttt{nonzero}. Đồng thời, nếu phép cộng có tràn nhớ lên bit~30 và bit~0 bằng~1
  ($\texttt{lost\_carry} = M[30]\land M[0]$), S4 jam ngay bit này vào vị trí bit~1 của
  $\texttt{norm\_seed[29:0]} = \{M[29:2], M[1]\lor\texttt{lost\_carry}, M[0]\}$ và OR
  vào \texttt{sticky\_raw} $= \texttt{s3\_tail\_q} \lor \texttt{lost\_carry}$.
\end{itemize}

\rtlfig{Datapath chi tiết tầng S3 (\sig{posit\_add\_arith\_comb}) và tầng S4 (\sig{posit\_add\_lod\_comb}): cộng/trừ độ lớn 31 bit, tiền tăng scale, dò LOD 31 bit và bảo toàn phần dư \sig{lost\_carry}.}{fig:adder-tail}{
\node[pin] (in1) at (-4.95,.76) {\sig{s2\_x\_frac\_q[25:0]}, \sig{s2\_y\_q[30:0]}, \sig{s2\_sub\_q}};
\node[pin] (in2) at (.20,.76) {\sig{s2\_sf\_q[10:0]}};
\node[pin] (in3) at (5.30,.76) {\sig{s2\_tail\_q}, dấu, cờ \sig{Z/N}, \sig{flags\_q}};

\node[unit,text width=41mm] (alu) at (-4.95,-1.05) {\hw{S3: cộng/trừ độ lớn 31 bit}\\
 $X=\{0,1,f_X,3'\mathrm{b0}\}$\\
 $M=\sig{subtract}\text{ ? }(X-Y):(X+Y)$};
\node[unit,text width=44mm] (inc) at (.20,-1.05) {\hw{S3: tiền tăng scale song song}\\
 $\sig{sf\_inc}[10:0]=\sig{sf\_base}+11'\mathrm{sd1}$};
\node[unit,text width=36mm] (m3) at (5.30,-1.05) {\hw{S3: giữ tail \& metadata}\\
 \sig{tail}, \sig{sign}, \sig{both\_zero},\\ \sig{is\_nar}, \sig{flags\_lo[3:0]}};

\draw[wire] (in1)--(alu);
\draw[wire] (in2)--(inc);
\draw[wire] (in3)--(m3);

\node[ff,text width=140mm] (r3) at (.15,-2.72) {\hw{S3}\enspace
 \sig{s3\_mag\_q[30:0]} \enspace|\enspace
 \sig{s3\_sf\_inc\_q[10:0]} \enspace|\enspace
 \sig{s3\_tail\_q}, \sig{s3\_sign\_q}, \sig{s3\_both\_zero\_q}, \sig{s3\_nar\_q}, \sig{s3\_flags\_q[3:0]}};

\draw[wire] (alu.south)--(r3.north -| alu);
\draw[wire] (inc.south)--(r3.north -| inc);
\draw[meta] (m3.south)--(r3.north -| m3);

\node[unit,text width=40mm] (lod) at (-4.95,-5.12) {\hw{S4: bộ dò LOD 31 bit}\\
 \sig{lod\_lzd\_core\_fixed2}\\
 $M[30:0]\to K[4:0],\ \sig{nonzero}$};
\node[unit,text width=46mm] (seed) at (.20,-5.12) {\hw{S4: seed 30 bit \& sticky\_raw}\\
 $\sig{lost\_carry}=M[30]\land M[0]$\\
 $\sig{seed}=\{M[29:2], M[1]{\lor}L, M[0]\}$\\
 $\sig{sticky\_raw}=\sig{tail}\lor\sig{lost\_carry}$};
\node[unit,text width=35mm] (m4) at (5.30,-5.12) {\hw{S4: giữ scale \& cờ}\\
 \sig{sf\_inc[10:0]}, \sig{sign},\\ \sig{both\_zero}, \sig{is\_nar}, \sig{flags\_lo}};

\draw[wire] (r3.south -| lod)--(lod.north);
\draw[wire] (r3.south -| seed)--(seed.north);
\draw[meta] (r3.south -| m4)--(m4.north);

\node[tapdot] at (-4.95,-3.65) {};
\draw[wire] (-4.95,-3.65) -- (-1.45,-3.65) -- (seed.north -| -1.45,0);
\node[tag] at (-3.15,-3.65) {$M[30:0]$};

\node[tapdot] at (5.30,-3.65) {};
\draw[meta] (5.30,-3.65) -- (1.85,-3.65) -- (seed.north -| 1.85,0);
\node[tag] at (3.65,-3.65) {\sig{s3\_tail\_q}};

\node[ff,text width=140mm] (r4) at (.15,-7.05) {\hw{S4}\enspace
 \sig{s4\_count\_q[4:0]}, \sig{s4\_nz\_q} \enspace|\enspace
 \sig{s4\_seed\_q[29:0]}, \sig{s4\_sticky\_q}\\
 \sig{s4\_sf\_inc\_q[10:0]} \enspace|\enspace
 \sig{s4\_sign\_q}, \sig{s4\_both\_zero\_q}, \sig{s4\_nar\_q}, \sig{s4\_flags\_q[3:0]}};

\draw[wire] (lod.south)--(r4.north -| lod);
\draw[wire] (seed.south)--(r4.north -| seed);
\draw[meta] (m4.south)--(r4.north -| m4);

\node[pin] (o) at (.15,-8.18) {Tới S5 (\sig{posit\_add\_norm\_comb}): dịch trái \sig{seed} theo $K$, tính $\sig{sf\_inc}-K$ và phân loại đầu ra};
\draw[wire] (r4)--(o);
}


Một trường hợp dễ kiểm là $\widehat P$ và C có cùng magnitude,
trái dấu và $\Delta sf=0$. Hai significand thẳng hàng,
hiệu bằng 0, LOD không tìm thấy bit 1 và Packer chọn mã Zero.
Nếu số hạng thứ hai chỉ nhỏ hơn một bit thấp, LOD thấy bit 1
ở xa đầu, scale giảm và bit thấp được kéo lên. Vì vậy một
thay đổi nhỏ trước adder có thể làm sai số tương đối tăng
đáng kể trong nhóm gần triệt tiêu. Đó là lý do Chương~\ref{ch:measurements}
dùng sai số tuyệt đối cho nhóm này.

\subsection{Tầng S5: Dịch trái chuẩn hóa, cập nhật scale và phân loại đầu ra}

Khối tổ hợp S5 (\texttt{posit\_add\_norm\_comb}, Hình~\ref{fig:adder-s5}) nhận
bộ đếm \texttt{leading\_count[4:0]} để dịch trái từ \texttt{norm\_seed[29:0]} qua
bộ dịch \texttt{dyn\_left\_shifter\_fixed6}. Vì phép dịch trái 30 bit tự đẩy bit ẩn~1 lên
vượt khỏi biên bit~28 của kết quả trong mọi chế độ (tràn nhớ $K=0$, bình thường
$K=1$, và triệt tiêu sâu $K\ge 2$), một phép trừ duy nhất
$\texttt{sf} = \texttt{sf\_inc} - K$ xử lý thống nhất cả ba chế độ mà không cần bộ
chọn tăng/giảm scale riêng.

Tại thanh ghi đầu ra S5, RTL chỉ chốt 29 bit fraction cao vào \texttt{frac\_hi\_q[28:0]}
và nối thêm 24 bit~0 hằng tại biên giao diện (\texttt{frac = \{frac\_hi\_q, 24'b0\}}),
nhờ vậy loại bỏ được 24 flip-flop luôn mang giá trị 0. Đồng thời, logic phân loại
\texttt{Zero/NaR} đặt tại tầng Packer thay vì chèn 41 MUX xóa bus số học trên đường
găng S5; tầng S5 chỉ cập nhật trực tiếp cờ ngoại lệ (\texttt{flags\_out[4]=is\_nar\_in},
và bit \texttt{inexact} chỉ OR thêm \texttt{sticky\_raw} khi kết quả khác Zero).

\rtlfig{Datapath chi tiết tầng S5 (\sig{posit\_add\_norm\_comb}): dịch trái chuẩn hóa 30 bit, trừ scale một bước và phân loại ngoại lệ không xóa bus dữ liệu.}{fig:adder-s5}{
\node[pin] (in1) at (-4.95,.82) {\sig{s4\_seed\_q[29:0]}, \sig{s4\_count\_q[4:0]}};
\node[pin] (in2) at (0,.82) {\sig{s4\_sf\_inc\_q[10:0]}};
\node[pin] (in3) at (4.95,.82) {\sig{s4\_nz\_q}, \sig{sticky}, dấu, \sig{Z/N}, \sig{flags\_lo}};

\node[unit,text width=41mm] (lsh) at (-4.95,-1.65) {\hw{Dịch trái chuẩn hóa 30 bit}\\
 \sig{dyn\_left\_shifter\_fixed6}\\
 $\{\sig{shifted\_hi[28:0]},\sig{lsb}\}=\sig{seed}\ll K$\\
 $\sig{frac\_hi}=\sig{shifted\_hi}$ (không xóa 0)};
\node[unit,text width=41mm] (sf) at (0,-1.65) {\hw{Cập nhật scale một bước}\\
 $\sig{count\_ext}=\{6'\mathrm{b0},K[4:0]\}$\\
 $\sig{sf[10:0]}=\sig{sf\_inc}-\sig{count\_ext}$\\
 gộp tràn $K{=}0$ và triệt tiêu $K{\ge}2$};
\node[unit,text width=43mm] (cls) at (4.95,-1.65) {\hw{Phân loại Zero/NaR \& cờ}\\
 $\sig{fz}=\neg\sig{nar\_in}\land(\sig{both\_zero}\lor\neg\sig{nz})$\\
 $\sig{sign}=\sig{nar\_in}\text{ ? }1:(\sig{fz}\text{ ? }0:\sig{sign\_base})$\\
 $\sig{inexact}|=\neg\sig{fz}\land\sig{sticky\_raw}$};

\draw[wire] (in1)--(lsh);
\draw[wire] (in2)--(sf);
\draw[wire] (in3)--(cls);
\node[tapdot] at (-4.95,.14) {};
\draw[wire] (-4.95,.14) -- (-1.35,.14) -- (sf.north -| -1.35,0);
\node[tag] at (-2.95,.14) {$K[4:0]$};

\node[ff,text width=142mm] (r5) at (0,-3.92) {\hw{S5}\enspace
 \sig{frac\_hi\_q[28:0]} \enspace|\enspace
 \sig{sf[10:0]} \enspace|\enspace
 \sig{sticky} \enspace|\enspace
 \sig{sign}, \sig{is\_zero}, \sig{is\_nar}, \sig{flags\_out[4:0]}, \sig{out\_valid}};

\draw[wire] (lsh.south)--(r5.north -| lsh);
\draw[wire] (sf.south)--(r5.north -| sf);
\draw[meta] (cls.south)--(r5.north -| cls);

\node[unit,text width=106mm] (pad) at (0,-5.32) {\hw{Khớp giao diện Packer (tiết kiệm 24 Flip-Flop hằng 0)}\\ $\sig{frac[52:0]}=\{\sig{frac\_hi\_q[28:0]},\ 24'\mathrm{b0}\}$};
\draw[wire] (r5.south)--(pad.north);

\node[pin] (o) at (0,-6.50) {Tới Packer P1 (\sig{posit\_pack\_prepare}):\\
 \sig{frac[52:0]}, \sig{sf[10:0]}, \sig{sticky}, \sig{sign}, \sig{is\_zero}, \sig{is\_nar}, \sig{flags\_out}};
\draw[wire] (pad.south)--(o.north);
}


\subsection{Truy vết thanh ghi bộ cộng qua S1--S5}

Để đối chiếu trực tiếp với hai phép tính tay tại Chương~\ref{ch:arithmetic}
($3{,}5 + 0{,}75 = 4{,}25$ và $3{,}5 + (-3{,}25) = 0{,}25$),
Bảng~\ref{tab:adder-trace-examples} ghi lại giá trị chốt trong các thanh ghi
của \texttt{posit\_add} tại ranh giới từng tầng \texttt{S1..S5}. Cùng một bộ
dò LOD 31 bit (\texttt{K=0} khi có tràn nhớ tại bit~30; \texttt{K=4} khi
bit dẫn đầu rơi xuống bit~26 do triệt tiêu) và cùng một phép trừ
\texttt{sf = sf\_inc - K} xử lý thống nhất cả hai chế độ mà không cần bộ
phân nhánh tăng/giảm scale riêng.

\begin{engineeringtable}{Truy vết giá trị thanh ghi qua năm tầng của \texttt{posit\_add} cho trường hợp tràn nhớ ($4{,}25$) và triệt tiêu ($0{,}25$).}{tab:adder-trace-examples}
\begin{tabularx}{\linewidth}{@{}>{\raggedright\arraybackslash}p{16mm} >{\raggedright\arraybackslash}p{62mm} >{\raggedright\arraybackslash}p{62mm}@{}}
\toprule
Tầng & Ca 1: $3{,}5 + 0{,}75 = 4{,}25$ (cộng tràn nhớ) & Ca 2: $3{,}5 + (-3{,}25) = 0{,}25$ (trừ triệt tiêu) \\
\midrule
Đầu vào &
$A{:}\ sf{=}1,\ f{=}\texttt{26'h3000000},\ s{=}0$\newline
$B{:}\ sf{=}-1,\ f{=}\texttt{26'h2000000},\ s{=}0$ &
$A{:}\ sf{=}1,\ f{=}\texttt{26'h3000000},\ s{=}0$\newline
$B{:}\ sf{=}1,\ f{=}\texttt{26'h2800000},\ s{=}1$ \\
\addlinespace[3pt]
\texttt{S1} &
$\sig{choose\_a}{=}1$, $\sig{s1\_shift\_q}{=}2$, $\sig{s1\_sub\_q}{=}0$, $\sig{s1\_sf\_q}{=}1$ &
$\sig{choose\_a}{=}1$, $\sig{s1\_shift\_q}{=}0$, $\sig{s1\_sub\_q}{=}1$, $\sig{s1\_sf\_q}{=}1$ \\
\texttt{S2} &
$X{=}\texttt{31'h38000000}$, $\sig{s2\_y\_q}{=}\texttt{31'h0c000000}$, $\sig{tail}{=}0$ &
$X{=}\texttt{31'h38000000}$, $\sig{s2\_y\_q}{=}\texttt{31'h34000000}$, $\sig{tail}{=}0$ \\
\texttt{S3} &
$\sig{s3\_mag\_q}{=}\texttt{31'h44000000}$ (bit~30 $=1$), $\sig{s3\_sf\_inc\_q}{=}2$ &
$\sig{s3\_mag\_q}{=}\texttt{31'h04000000}$ (bit~26 $=1$), $\sig{s3\_sf\_inc\_q}{=}2$ \\
\texttt{S4} &
$\sig{s4\_count\_q}{=}0$, $\sig{s4\_seed\_q}{=}\texttt{30'h04000000}$, $\sig{sticky}{=}0$ &
$\sig{s4\_count\_q}{=}4$, $\sig{s4\_seed\_q}{=}\texttt{30'h04000000}$, $\sig{sticky}{=}0$ \\
\texttt{S5} &
$\sig{sf}{=}2-0=2$, $\sig{frac\_hi\_q}{=}\texttt{29'h02000000}$ ($1{,}0625\times 2^2=4{,}25$) &
$\sig{sf}{=}2-4=-2$, $\sig{frac\_hi\_q}{=}\texttt{29'h00000000}$ ($1{,}0\times 2^{-2}=0{,}25$) \\
\bottomrule
\end{tabularx}
\end{engineeringtable}

\subsection{Packer P1 tạo G/R/S và phân loại}

Quá trình đóng gói hai tầng của bộ cộng/tích được mô tả trong
Algorithm~\ref{alg:posit-pack}, tương ứng với hai mô-đun
\texttt{posit\_pack\_prepare} (P1) và \texttt{posit\_pack\_finish} (P2).

\begin{reportpseudocode}{alg:posit-pack}{Đóng gói và làm tròn RNE hai tầng (Two-Stage Posit(32,3) RNE Packing)}
  {Kết quả chuẩn hóa $(s, sf[10:0], f[52:0], \mathrm{sticky}, z, \eta, \mathrm{flags_{in}}[4:0])$.}
  {Từ mã đóng gói Posit(32,3) $d[31:0]$ và cờ trạng thái $\mathrm{flags}[4:0]$.}
  \begin{tabularx}{\linewidth}{@{}r >{\raggedright\arraybackslash}X >{\raggedright\arraybackslash}p{.33\linewidth}@{}}
    \algline{1}{\textbf{P1} \quad $r_0\gets \neg sf[10]$, $\mathrm{offset}[4:0]\gets sf[7:3]\oplus\{5\{sf[10]\}\}$}{Xác định bit regime đầu và độ dịch không rẽ nhánh.}
    \algline{2}{$\mathrm{payload}[58:0]\gets \{r_0, \neg r_0, sf[2:0], f[52:0], 1'b0\}$}{Ghép 59 bit gồm regime, bit dừng, mũ và phân số.}
    \algline{3}{$\mathrm{shifted}\gets (\mathrm{payload}\gg_r \mathrm{offset})$}{Dịch phải điền $r_0$ qua 5 tầng MUX để tạo regime.}
    \algline{4}{$m_{\mathrm{trunc}}[30:0]\gets \mathrm{shifted}[58:28]$, $G\gets \mathrm{shifted}[27]$, $R\gets \mathrm{shifted}[26]$}{Tách độ lớn 31 bit và hai bit làm tròn Guard, Round.}
    \algline{5}{$S\gets \mathrm{sticky}\lor (\bigvee f[24:0])\lor \bigvee_{k=0}^{30}(\{sf[2:0],f[52:25]\}_k \land (k<\mathrm{offset}))$}{Tính bit Sticky song song bằng mặt nạ nhiệt kế.}
    \algline{6}{$\mathrm{sel}\gets \mathrm{Classify}(\eta, z, sf, f, \mathrm{sticky})\in\{\mathrm{NAR},\mathrm{ZERO},\mathrm{MAX},\mathrm{MIN},\mathrm{NORM}\}$}{Ưu tiên NaR, Zero và bão hòa tại biên $sf\in[-240,240]$.}
    \algline{7}{\textbf{P2} \quad $\mathrm{round\_up}\gets G\land (m_{\mathrm{trunc}}[0]\lor R\lor S)$}{Kiểm tra điều kiện làm tròn gần nhất chẵn (RNE).}
    \algline{8}{$\mathrm{base}\gets \{1'b0, m_{\mathrm{trunc}}\}\oplus \{32\{s\}\}$, $\mathrm{inc}\gets s\oplus \mathrm{round\_up}$}{Gộp bù hai khi $s=1$ và lượng tăng RNE thành một bộ cộng.}
    \algline{9}{$d_{\mathrm{round}}\gets \mathrm{base}+\mathrm{inc}$}{Tính từ mã đã làm tròn trên nhánh bình thường.}
    \algline{10}{\textbf{if} $\mathrm{sel}=\mathrm{NAR}$ \textbf{then} $d\gets \texttt{32'h80000000}$, $\mathrm{flags}\gets \texttt{5'b10000}$}{Xuất mẫu NaR và đặt cờ ngoại lệ NaR.}
    \algline{11}{\textbf{else if} $\mathrm{sel}\in\{\mathrm{ZERO},\mathrm{MAX},\mathrm{MIN}\}$ \textbf{then} $d\gets \mathrm{Const}(\mathrm{sel},s)$}{Xuất $0$, $\pm\mathrm{maxpos}$ hoặc $\pm\mathrm{minpos}$ theo $\mathrm{sel}$.}
    \algline{12}{\textbf{else} $d\gets \mathrm{Clamp}(m_{\mathrm{trunc}}, s, d_{\mathrm{round}})$, $\mathrm{flags}[1]\gets \mathrm{flags_{in}}[1]\lor G\lor R\lor S$}{Chặn tràn biên max/min và cập nhật cờ \texttt{inexact}.}
    \algline{13}{\textbf{return} $d, \mathrm{flags}$}{Chốt kết quả Posit(32,3) cuối đường ống P2.}
  \end{tabularx}
\end{reportpseudocode}

P1 nhận scale signed 11 bit, fraction 53 bit, dấu, sticky và
metadata ngoại lệ. Bit regime đầu là nghịch đảo dấu của scale.
Với $k\ge0$, offset bằng phần thấp của $k$. Với $k<0$, offset
là $|k|-1$ trên miền normal. Payload 59 bit ghép bit regime đầu,
terminator, exponent 3 bit, fraction 53 bit và một bit 0.
Bộ dịch phải điền bit regime đầu vào các vị trí trống bên trái,
tạo magnitude giữ lại từ \texttt{shifted[58:28]}. Hai bit
tiếp theo là guard và round.

\rtlfig{Packer cuối: P1 dựng magnitude và G/R/S song song với phân loại; P2 ghép RNE với bù hai và chọn mã. Mọi đường đều đi qua hai tầng valid/ready.}{fig:packer-stages}{
\node[pin] (i) at (0,.69) {\sig{sf[10:0]}, \sig{frac[52:0]}, sticky, dấu, Zero/NaR, flags 5};
\node[unit,text width=42mm] (asm) at (-5,-1.46) {\hw{Regime + payload}\\$r=\neg sf[10]$\\$o=sf[7:3]\oplus\{5\{sf[10]\}\}$\\$P=\{r,\neg r,sf[2:0],f,0\}$};
\node[unit,text width=42mm] (mask) at (0,-1.46) {\hw{Mặt nạ sticky 31 bit}\\$U=\{sf[2:0],f[52:25]\}$\\chọn $U_j$ nếu $j<o$\\$S=\mathrm{sticky}\lor |f[24:0]\lor |\mathrm{mask}$};
\node[unit,text width=42mm] (class) at (5,-1.46) {\hw{Phân loại song song}\\NaR $\to$ Zero $\to$ MAX $\to$ MIN\\MAX: $sf\ge240$; MIN: $sf<-240$\\\sig{special\_sel}, flags};
\foreach \n in {asm,mask,class}{\draw[wire] (i.south)--++(0,-.24)-|(\n.north);}
\node[unit,text width=42mm] (sh) at (-5,-3.91) {\hw{Dịch phải 59 bit}\\$W=P\gg o$; điền $r$\\$m=W[58:28]$\\$G=W[27]$, $R=W[26]$};
\draw[wire] (asm)--node[tag,right]{59, $o$:5, $r$:1}(sh);
\draw[wire] (asm.east)--(mask.west);
\node[ff,text width=145mm] (r1) at (0,-5.6) {\hw{P1}\enspace $m[30:0]$\enspace|\enspace $G,R,S$\enspace|\enspace dấu\enspace|\enspace \sig{special\_sel[2:0]}, flags};
\draw[wire] (sh)--(r1.north -| sh);
\draw[wire] (mask)--node[tag,right]{$S$:1}(r1.north -| mask);
\draw[meta] (class)--node[tag,right]{loại, flags}(r1.north -| class);
\node[unit,text width=66mm] (r) at (-3.73,-7.34) {\hw{RNE + mã dấu: một phép cộng}\\
 $u=G\land(m[0]\lor R\lor S)$\\
 $D_{\mathrm{r}}=(\{0,m\}\oplus\{32\{s\}\})+(s\oplus u)$};
\node[unit,text width=62mm] (end) at (3.73,-7.34) {\hw{Biên và cờ}\\
 test $m=0$ / $m=\mathtt{7fff\,ffff}$\\
 chọn min/max theo dấu\\nhánh thường: inexact OR G/R/S};
\draw[wire] (r1.south -| r)--(r);
\draw[wire] (r1.south -| end)--(end);
\node[pick,text width=81mm] (mux) at (0,-9.38) {\hw{Chọn mã theo}~\sig{special\_sel}~\hw{và biên}};
\draw[wire] (r.south)--++(0,-.32)-|(mux.north -| r.east);
\draw[wire] (end.south)--++(0,-.32)-|(mux.north -| end.west);
\node[ff,text width=95mm] (r2) at (0,-10.65) {\hw{P2} \quad \sig{d[31:0]}, \sig{flags[4:0]}, \sig{out\_valid}};
\draw[wire] (mux)--(r2);
}



P1 tính sticky song song với bộ dịch. Fraction thấp
\texttt{frac[24:0]} luôn nằm dưới round, nên OR trực tiếp với
sticky đầu vào. Các bit cao hơn chỉ tham gia sticky khi vị
trí của chúng nhỏ hơn offset tương ứng. Mỗi bit dùng một so
sánh hằng theo offset và cổng AND với dữ liệu. Cuối cùng OR
các bit được chọn. Nhờ vậy phép tìm sticky không đi nối tiếp
qua toàn bộ đường barrel shifter.

\subsection{Packer P2 làm tròn, kẹp và mã hóa dấu}

Với magnitude giữ $m$ và quyết định tăng $r\in\{0,1\}$,
đầu ra dương bình thường là $m+r$. Với số âm, biểu thức
$-(m+r)=\mathop{\mathrm{NOT}}m+(1-r)$ trên 32 bit cho phép
ghép việc tăng RNE và mã hóa bù hai. RTL chọn XOR theo dấu
và dùng increment \texttt{sign XOR round\_up}. Đường chọn
ngoại lệ vẫn có thể trả Zero, NaR, minpos hoặc maxpos.

P1 mang theo chỉ dấu nhánh đặc biệt và flags. NaR ưu tiên,
sau đó Zero và các nhánh biên. Với nhánh thường,
\texttt{inexact} OR thông tin mất từ trước với G/R/S.
\texttt{approx\_cut} vẫn chỉ báo còn số hạng tại điểm dừng
bộ nhân. Việc tách hai cờ giúp scoreboard phân biệt
lượng tử hóa ở Packer với giới hạn vòng.

\section{Phạm vi đối chứng NF và quire}

Các khối P1/P2, nhân lặp, adder và Packer cuối của MAC-NF và MAC-F được
triển khai đồng bộ trên cùng bộ nguồn RTL dùng chung khối nhân. NF có đường tích đi qua
Packer tích hai tầng P1/P2 (\texttt{posit\_product\_pack}) và Parser tích
(\texttt{posit\_magnitude\_parse}) rồi đến adder (Hình~\ref{fig:nf-f-split}).
Cụ thể, ở tầng P1, mạch tổ hợp \texttt{posit\_pack\_prepare} dịch phải 59 bit
và tạo mặt nạ nhiệt sticky rồi chốt \texttt{p1\_mag\_trunc\_q[30:0]},
\texttt{G/R/S}, \texttt{p1\_sign\_q}, \texttt{p1\_special\_sel\_q[2:0]},
\texttt{p1\_flags\_q[4:0]} và \texttt{p1\_valid\_q}. Ở tầng P2, mạch tổ hợp
\texttt{posit\_pack\_finish} làm tròn RNE 32 bit (\texttt{sign=0}) và ghép nhánh
ngoại lệ rồi chốt \texttt{d\_q[31:0]}, \texttt{flags\_q[4:0]}, \texttt{sign\_q}
và \texttt{p2\_valid\_q}. Tiếp đó, \texttt{posit\_magnitude\_parse} giải mã tổ hợp
trước khi bank tích NF chốt \texttt{norm\_valid\_q} để đưa vào S1.

\rtlfig{Biên tích trước S1 của MAC-F và MAC-NF. Hai cột biểu diễn hai thiết kế riêng; cả hai giữ C từ lúc Parser retire, còn \sig{posit\_product\_pack} gồm hai tầng pipeline P1/P2 rõ mạch tổ hợp và FF.}{fig:nf-f-split}{
\node[unit,text width=116mm] (src) at (0,0) {\hw{Adapter sau nhân:}~\sig{mul\_prepack\_norm\_comb}\\
 \sig{acc[25]} chọn $(sf+1,\{\sig{acc[24:0]},0\})$ hoặc $(sf,\{\sig{acc[23:0]},00\})$};
\node[head] (hf) at (-3.5,-1.18) {MAC-F};
\node[head] (hn) at (3.65,-1.18) {MAC-NF};
\node[unit,text width=66mm] (p1c) at (3.65,-2.52) {\hw{P1 tổ hợp:}~\sig{posit\_pack\_prepare}\\
 dịch phải 59 bit (\sig{frac}=\{26, 27'b0\}, \sig{sticky}=0)\\mặt nạ nhiệt sticky; phân loại \sig{special\_sel}};
\node[ff,text width=66mm] (p1r) at (3.65,-4.48) {\hw{P1 FF:}~\sig{posit\_product\_pack} (\sig{p1\_ready})\\
 \sig{p1\_mag\_trunc\_q[30:0]}, \sig{G/R/S}, \sig{p1\_sign\_q}\\\sig{p1\_special\_sel\_q[2:0]}, \sig{p1\_flags\_q[4:0]}, \sig{p1\_valid\_q}};
\node[unit,text width=66mm] (p2c) at (3.65,-6.42) {\hw{P2 tổ hợp:}~\sig{posit\_pack\_finish} (\sig{sign=0})\\
 \sig{round\_up} RNE 32 bit, MUX \sig{special\_sel}/biên, cờ \sig{inexact}};
\node[ff,text width=66mm] (p2r) at (3.65,-8.18) {\hw{P2 FF:}~\sig{posit\_product\_pack} (\sig{p2\_ready})\\
 \sig{d\_q[31:0]} (magnitude), \sig{flags\_q[4:0]}, \sig{sign\_q}, \sig{p2\_valid\_q}};
\node[unit,text width=66mm] (mp) at (3.65,-10.02) {\sig{posit\_magnitude\_parse}\\XOR + sentinel $\to$ LOD 31\\tái tạo frac 26, $sf$ 11, $s,Z,N$};
\node[ff,text width=66mm] (nr) at (3.65,-11.96) {\hw{Bank tích NF} (\sig{product\_capture})\\\sig{norm\_valid\_q}, \sig{product\_frac\_q}, \sig{norm\_sf\_q}, $s/Z/N$, \sig{norm\_flags\_q}};
\node[unit,text width=54mm] (af) at (-3.5,-13.62) {\hw{S1 nhận tích F và C}};
\node[unit,text width=66mm] (an) at (3.65,-13.62) {\hw{S1 nhận tích NF và C}};
\draw[wire] (src.south) -- ++(0,-.26) -| (hn);
\draw[wire] (src.south) -- ++(0,-.26) -| (hf);
\draw[wire] (hn)--node[tag,right]{\sig{core\_done}}(p1c);
\draw[wire] (p1c)--(p1r);
\draw[wire] (p1r)--(p2c);
\draw[wire] (p2c)--(p2r);
\draw[wire] (p2r)--node[tag,right]{magnitude 32 + \sig{sign\_q}}(mp);
\draw[wire] (mp)--node[tag,right]{\sig{product\_capture}}(nr);
\draw[wire] (hf)--(af) node[pos=.28,tag,right]{frac 26 và $sf$ 11\\đưa thẳng vào S1\\với \sig{core\_done}};
\draw[wire] (nr)--(an);
\draw[meta] (p2r.west)--++(-.68,0) |- (nr.west) node[pos=.42,tag,left]{\sig{flags\_q[4:0]}};
}
 Khi giải thích
chênh lệch tài nguyên, so hai kết quả trong cùng một lượt
tổng hợp. Khi giải thích chu kỳ, dùng cùng dãy A/B và cùng
điều kiện backpressure; không lấy độ sâu hình khối để suy
ra tần số hay throughput.

Với quy ước đo tuần tự từng phép một (mỗi lần 1 phép đang chạy, \texttt{out\_ready=1}, đếm gộp từ cạnh chấp nhận đầu vào đến khi hoàn tất rút kết quả và \texttt{in\_ready} sẵn sàng ở cạnh kế tiếp), số chu kỳ của một phép MAC khác 0 phụ thuộc số bit~1 được chọn $t = \mathrm{popcount}(X_{\mathrm{mant}})$ ở bộ nhân như sau:
\begin{itemize}
\item \textbf{MAC-F:} $T_F(t) = 13 + \max(1,t) \in [14, 25]$ chu kỳ với $0 \le t \le 12$ (Parser \texttt{P1/P2}: 2; nạp ngữ cảnh/\texttt{T1}: 3; lặp \texttt{T2/T3}: $\max(1,t)$; Adder \texttt{S1..S5}: 5; Packer \texttt{P1/P2}: 2; chuyển tiếp retire sang \texttt{in\_ready}: 1; riêng khoảng cách từ cạnh nhận đến \texttt{out\_valid} là $12 + \max(1,t)$ cạnh).
\item \textbf{MAC-NF:} $T_{NF}(t) = 16 + \max(1,t) = T_F(t) + 3 \in [17, 28]$ chu kỳ. Do \texttt{posit\_magnitude\_parse} là mạch tổ hợp, phần chênh đúng $+3$ chu kỳ đến từ 2 tầng FF \texttt{P1/P2} của \texttt{posit\_product\_pack} và 1 tầng FF \texttt{norm\_valid\_q} tại bank tích NF trước tầng cộng.
\item \textbf{MAC-Quire:} Lệnh \texttt{CLEAR} mất $3$ chu kỳ; lệnh \texttt{MAC} với tích khác 0/NaR mất $36 + \max(1,t) = T_F(t) + 23 \in [37, 48]$ chu kỳ do FSM đi qua 30 chu kỳ xử lý tại bank tích lũy 27 từ sau \texttt{core\_done} (thay vì 7 chu kỳ nối thẳng qua Adder/Packer ở MAC-F), còn tích bằng 0 hoặc NaR bỏ qua vòng bank 27 bước để xong sau 1 chu kỳ; lệnh \texttt{READ} mất $59$ chu kỳ để quét/giảm 27 từ quire và đóng gói \texttt{posit\_pack}.
\end{itemize}

Quire dùng mô-đun đỉnh \texttt{posit\_mac\_quire} và các module backend
cập nhật theo từng từ 36 bit. Lệnh MAC căn tích
sang lưới $2^{-480}$ và giữ $Q$ có dấu 972 bit. CLEAR đưa
$Q$/count/status về trạng thái khởi tạo; READ quét $Q$ để
tạo fraction và sticky cho Packer. Chi tiết hiệu quả tối ưu hóa của kiến trúc
backend theo word (từ phiên bản dịch toàn vector sang cập nhật theo từ 36 bit)
và số chu kỳ đo trên chuỗi tích lũy được trình bày cụ thể trong Chương~\ref{ch:measurements}.

\rtlfig{Backend quire cập nhật theo từ 36 bit. MAC cập nhật từng word rồi commit Q; READ đọc Q đã commit, dò bit dẫn đầu theo word và làm tròn đúng một lần.}{fig:quire-context}{
\node[head] at (-3.85,0) {Ghi: lệnh MAC};
\node[head] at (3.85,0) {Đọc: lệnh READ};
\node[ff,text width=63mm] (prod) at (-3.85,-1.15) {\hw{Chốt tích từ core}\\magnitude 26 Q24, $sf$ 11, dấu/cờ};
\node[unit,text width=63mm] (al) at (-3.85,-3.24) {\hw{Căn lưới $2^{-480}$}\\
 $h=sf+456$; hệ số signed 27 bit\\
  index $=h_+\ \mathrm{div}\ 36$; offset $=h_+\bmod36$\\
  $h_+=\max(h,0)$ trên nhánh hợp lệ};
\draw[wire] (prod)--(al);
\node[unit,text width=63mm] (w) at (-3.85,-5.53) {\hw{Chọn term của word đang xử lý}\\
 dịch hệ số trong 63 bit; word 36 bit\\
 chọn 0, lát thấp, lát cao hoặc dấu};
\draw[wire] (al)--node[tag,right]{hệ số 27, index 5, offset 6}(w);
\node[unit,text width=63mm] (sum) at (-3.85,-7.55) {\hw{Cộng theo word 0 $\to$ 26}\\
 \sig{q\_words[step]} + term + carry\\
 tổng 37 bit; bit 36 truyền vòng sau};
\draw[wire] (w)--(sum);
\node[ff,text width=63mm] (pend) at (-3.85,-9.43) {\hw{Bank pending}\\
 ghi word tổng; commit khi xong word 26 và không overflow};
\draw[wire] (sum)--node[tag,right]{tổng thấp 36}(pend);
\node[ff,text width=64mm] (q) at (3.85,-1.15) {\hw{Q đã commit}\\27 word $\times$ 36 bit; signed 972 bit};
\draw[wire] (pend.south)--(-3.85,-10.26)--(.28,-10.26) -- (.28,-1.15) -- (q.west);
\node[tag,rotate=90] at (.28,-9.15) {commit};
\draw[wire] ($(q.west)+(0,.32)$)--(-.24,-.83)|-(sum.east);
\node[unit,text width=63mm] (fetch) at (3.85,-3.02) {\hw{FETCH: lấy độ lớn word}\\
 chọn word theo step từ thấp lên cao\\
 XOR theo dấu Q, cộng carry $\to$ word 36 bit};
\draw[wire] (q)--(fetch);
\node[ff,text width=63mm] (abs) at (3.85,-4.78) {\hw{Chốt}~\sig{abs\_word\_q}~\hw{và carry}};
\draw[wire] (fetch)--(abs);
\node[unit,text width=63mm] (red) at (3.85,-6.54) {\hw{REDUCE: dò và tạo cửa sổ}\\
 ưu tiên bit 1 cao nhất trong word\\
 cửa sổ 3 word dịch $35-\mathrm{lead\_bit}$};
\draw[wire] (abs)--(red);
\draw[enable] (red.east)--(7.56,-6.54) |- (fetch.east);
\node[ff,text width=63mm] (out) at (3.85,-8.67) {\hw{Cập nhật khi word $\ne0$}\\
 $sf=36\,\mathrm{step}+\mathrm{lead\_bit}-480$\\
 frac 53; sticky OR phần thấp\\và các word cũ};
\draw[wire] (red)--(out);
\node[unit,text width=63mm] (p) at (3.85,-10.52) {\hw{Packer RNE sau word 26}};
\draw[wire] (out)--(p);
\node[pin] (r) at (3.85,-11.6) {\sig{d[31:0]}, flags, \sig{out\_cmd=READ}};
\draw[wire] (p)--(r);
\node[explain,text width=63mm] at (-3.85,-11.6) {CLEAR xóa Q và trạng thái.\\CLEAR/MAC trả ack với $d=0$.};
}



\subsection{Commit quire và phép đọc giữ trạng thái}

Một MAC hợp lệ tạo term trên lưới quire từ dấu, magnitude
Q24 và scale. Phần căn tích chuẩn bị descriptor gồm hệ số,
chỉ số word và offset. Backend cập nhật theo từng word với
carry/borrow qua ranh giới, rồi commit state hoàn chỉnh.
Trong thời gian cập nhật, giao diện giữ một lệnh outstanding.
MAC chỉ trả acknowledge sau khi state đã cập nhật và count
tăng. Nếu sink kéo ready thấp, response giữ ổn định và không
commit lặp. C được giữ ở driver đối với NF/F; quire có state
trong backend và không có cổng C.

READ xác định dấu và bit dẫn đầu của Q, dựng một cửa sổ
fraction và tổng hợp sticky cho các bit thấp hơn cửa sổ.
Packer dùng dữ liệu đó để làm tròn đầu ra. READ giữ Q,
count và trạng thái loss của chuỗi. CLEAR tạo epoch mới
với Q và count bằng 0. Khi count đã đạt 1.024, lệnh MAC
tiếp bị chặn; READ và CLEAR vẫn có vai trò kết thúc hoặc
khởi tạo lại chuỗi. Điểm này được sử dụng trong benchmark:
mỗi chuỗi dài 1.024 có đúng 1.024 MAC rồi READ.

\begin{engineeringtable}{Giao diện lệnh của \texttt{posit\_mac\_quire} theo đặc tả kỹ thuật.}{tab:quire-command-flow}
\begin{tabularx}{\linewidth}{@{}>{\raggedright\arraybackslash}p{22mm}Y Y@{}}
\toprule
Lệnh & Điều kiện và tác động state & Response \\
\midrule
CLEAR & Xóa Q, count và trạng thái của epoch. & Ack với out\_cmd=CLEAR, d=0. \\
MAC & Nhận A/B và cấu hình khi còn dung lượng, tạo và commit term. & Ack với d=0, flags phản ánh trạng thái/loss hiện tại. \\
READ & Quét state Q đã commit, giữ state. & D là kết quả RNE và flags của READ. \\
\bottomrule
\end{tabularx}
\end{engineeringtable}

